% concatenated from QCLT-SWE-drift11.tex
\documentclass[12pt]{article}

%% BEGIN Setup

\usepackage{amsmath,amsthm,amsfonts,amssymb, color,xcolor,subcaption,graphicx,enumerate} %,leftindex}

%\usepackage{showkeys}


\newtheorem{theorem}{Theorem}[section]
\newtheorem{corollary}[theorem]{Corollary}
\newtheorem{example}[theorem]{Example}
\newtheorem{assumption}[theorem]{Assumption}
\newtheorem{hypothesis}[theorem]{Hypothesis}
\newtheorem{proposition}[theorem]{Proposition}
\newtheorem{lemma}[theorem]{Lemma}
\newtheorem{definition}[theorem]{Definition}
\newtheorem{remark}[theorem]{Remark}

\def\cA{\mathcal{A}}
\def\cB{\mathcal{B}}
\def\cC{\mathcal{C}}
\def\cD{\mathcal{D}}
\def\cE{\mathcal{E}}
\def\cF{\mathcal{F}}
\def\cG{\mathcal{G}}
\def\cH{\mathcal{H}}
\def\cI{\mathcal{I}}
\def\cJ{\mathcal{J}}
\def\cL{\mathcal{L}}
\def\cM{\mathcal{M}}
\def\cN{\mathcal{N}}
\def\cP{\mathcal{P}}
\def\cR{\mathcal{R}}
\def\cQ{\mathcal{Q}}
\def\cS{\mathcal{S}}
\def\cT{\mathcal{T}}
\def\cU{\mathcal{U}}
\def\cV{\mathcal{V}}
\def\cX{\mathcal{X}}
\def\cZ{\mathcal{Z}}

\def\bC{\mathbb{C}}
\def\bD{\mathbb{D}}
\def\bE{\mathbb{E}}
\def\bL{\mathbb{L}}
\def\bQ{\mathbb{Q}}
\def\bN{\mathbb{N}}
\def\bP{\mathbb{P}}
\def\bR{\mathbb{R}}
\def\bS{\mathbb{S}}
\def\bZ{\mathbb{Z}}

\newcommand{\fH}{\mathfrak{H}}

\def\wP{\widetilde{\cP}}
\def\Pb{\widetilde{\cP}_b}

\def\e{\varepsilon}
\def\k{\kappa}

\def\s{\sigma}

\def\SD{\mathbb{S}_{\mathbb{D}}}
\def\N{\overline{N}}
\def\barD{\overline{\mathbb{D}}_0}
\def\oZ{\overline{Z}}

\def\blue{\color{blue}}

\newcommand{\les}{\lesssim}

\newcommand{\Lip}{\operatorname{Lip}}

\def\bs{{\bf s}}
\def\bt{{\bf t}}
\def\bu{{\bf u}}
\def\bv{{\bf v}}
\def\bx{{\bf x}}
\def\by{{\bf y}}
\def\bz{{\bf z}}

\newcommand{\fm}{\mathfrak{m}}

\topmargin -0.4in
\headsep 0.4in
\textheight 9.0in
\oddsidemargin 0.02in
\evensidemargin 0.15in
\textwidth 6.3in

%% END Setup

\begin{document}

%% BEGIN Intro

\title{Gaussian fluctuations for the stochastic wave equation with drift}

\author{Raluca M. Balan\footnote{Corresponding author. University of Ottawa, Department of Mathematics and Statistics, 150 Louis Pasteur Private, Ottawa, Ontario, K1N 6N5, Canada. E-mail address: rbalan@uottawa.ca.} \footnote{Research supported by a grant from the Natural Sciences and Engineering Research Council of Canada.}
\and
William D. Stephenson \footnote{University of Ottawa, Department of Mathematics and Statistics, 150 Louis Pasteur Private, Ottawa, Ontario, K1N 6N5, Canada. E-mail address: wstep051@uottawa.ca.}
}

\date{September 17, 2026}
\maketitle

\begin{abstract}
\noindent In this article, we study Gaussian fluctuations of spatial averages of the solution to the stochastic wave equation with a nonlinear drift and multiplicative Gaussian noise in dimensions one and two. The noise is white in time, and its spatial covariance is either integrable or given by a Riesz kernel; space-time white noise in dimension one is also included. We establish spatial ergodicity and convergence of the rescaled covariance, and prove a quantitative central limit theorem for the centered spatial average over a ball of radius $R$. The bounds in total variation distance are of order $R^{-d/2}$ in the integrable case and $R^{-\beta/2}$ for a Riesz kernel of order $\beta$. We also obtain a functional central limit theorem in the space of continuous functions. The nonlinear drift creates additional difficulties in the estimation of the second Malliavin derivative and in the analysis of spatial averages. Our approach combines a second-order Gaussian Poincar\'e inequality with spatially integrated estimates for the second Malliavin derivative, exploiting the compact support of the wave kernel instead of a pointwise product-kernel bound. Further arguments based on the Clark-Ocone formula are used to control the drift contributions in the covariance and temporal-increment estimates.
\end{abstract}

\noindent {\em MSC 2020:} Primary 60H15; Secondary 60G60, 60G51
%60H15=SPDEs
%60G60=random fields
%60G51: Processes with independent increments; L\'evy processes

\vspace{1mm}

\noindent {\em Keywords:} stochastic wave equation, random fields, Malliavin calculus, Stein method, normal approximations

\pagebreak

\tableofcontents




\section{Introduction}

In this article, we consider the stochastic wave equation:
\begin{align}
\label{SWE}
	\begin{cases}
		\dfrac{\partial^2 u}{\partial t^2} (t,x)
		=  \Delta u (t,x)+\sigma\big(u(t,x)\big) \dot{W}(t,x)+b\big(u(t,x) \big), \ 
		t>0, \ x \in \bR^d \ (d\leq 2), \\
		u(0,x) = 1, \dfrac{\partial u}{\partial t} (0,x)=0, \quad x \in \bR^d,
	\end{cases}
\end{align}
driven by a spatially-homogeneous Gaussian noise which is white in time. More precisely, $W=\{W(\varphi);\varphi \in \cD(\bR_{+} \times \bR^d)\}$ is a zero-mean Gaussian process defined on a complete probability space $(\Omega,\cF,\bP)$, with covariance
\[
\bE[W(\varphi)W(\psi)]=\int_0^{\infty} \int_{\bR^d} \int_{\bR^d}\varphi(t,x)\psi(t,y)f(x-y)dxdydt=:\langle \varphi,\psi \rangle_{\cH},
\] 
where $f:\bR^d \to [0,\infty]$ is a tempered non-negative definite function.  By Bochner-Schwartz theorem, $f$ is the  Fourier transform of a tempered measure $\mu$ on $\bR^d$. Hence, for any $\varphi,\psi \in \cS(\bR^d)$,
\[
\int_{\bR^d} \int_{\bR^d}\varphi(x)\psi(y) f(x-y)dxdy=\int_{\bR^d}\cF \varphi(\xi) \overline{\cF \psi(\xi)}\mu(d\xi),
\]
where $\cS(\bR^d)$ is the space of rapidly decreasing functions on $\bR^d$, $\cF \varphi(\xi)=\int_{\bR^d} e^{-i\xi \cdot x} \varphi(x)dx$ is the Fourier transform of $\varphi$, and $x\cdot y$ denotes the usual inner product in $\bR^d$. Let $\cH$ denote the completion of $\cD(\bR_{+} \times \bR^d)$ with respect to the inner product $\langle \cdot,\cdot \rangle_{\cH}$. 

Throughout, we will assume that the functions $\sigma$ and $b$ are in $C^2(\bR)$ with $b',b'',\sigma',\sigma''$ bounded.

A predictable process $\{u(t,x);t\geq 0,x\in \bR^d\}$ is called a {\em (mild) solution} of \eqref{SWE}
if it satisfies the integral equation:
\begin{align}
\nonumber
u(t,x) & =1+\int_0^t \int_{\bR^d}G_{t-s}(x-y)\sigma\big(u(s,y)\big) W(ds,dy) \\
\label{mild soln}
& \quad +\int_0^t \int_{\bR^d} G_{t-s}(x-y) b\big(u(s,y)\big)dsdy,
\end{align}
where the stochastic integral is interpreted in the It\^o sense, and $G$ is the fundamental solution of the wave equation:
\begin{align*}
G_t(x)&=\frac{1}{2}\mathbf 1_{\{|x|<t\}} \quad \quad \quad \quad \quad \quad \mbox{if} \quad d=1\\
G_t(x)& =\frac{1}{2\pi} \frac{1}{\sqrt{t^2-|x|^2}} \mathbf 1_{\{|x|<t\}} \quad \mbox{if} \quad d=2.
\end{align*}
Here, $|\cdot|$ denotes the Euclidean norm. We will use the fact that for any $t>0$,
\begin{equation}
\label{int-G}
\int_{\bR^d}G_t(x)dx=t \quad \mbox{for} \quad d \in \{1,2\}.
\end{equation}




We recall that a random field $\{\Phi(t,x);t\geq 0,x\in \bR\}$ is {\em predictable} if it is measurable with respect to the predictable $\sigma$-field on $\Omega \times \bR_{+} \times \bR^d$, which is the minimal $\sigma$-field with respect to which all elementary processes are measurable. An {\em elementary process} is a linear combination of processes of the form
\[
\Phi(t,x)=Y \mathbf 1_{(a,b]}(t) \mathbf 1_{A}(x)
\]
where  $0\leq a<b$, $A \in \cB_b(\bR^d)$ and $Y$ is $\cF_a$-measurable. 




In \cite{dalang99}, it was proved that if $b$ and $\sigma$ are globally Lipschitz, and $\mu$ satisfies {\em Dalang's condition}:
\[
(D) \quad
\int_{\bR^d}\frac{1}{1+|\xi|^2}\mu(d\xi)<\infty,
\]
then equation \eqref{SWE} has a unique solution which satisfies: for any $p\geq 2$,
\begin{equation}
\label{def-KTp}
K_{T,p}:=\sup_{(t,x) \in [0,T] \times \bR^d}\bE|u(t,x)|^p<\infty \quad \mbox{for any} \quad T>0.
\end{equation}
In dimension $d=1$, Dalang's condition always holds. 

Since the noise is spatially homogeneous and the initial condition is constant, the process $\{u(t,x)\}_{x\in \bR^d}$ is strictly stationary.

For each fixed $t>0$, we consider the {\em spatial average} of the (centered) solution:
\[
F_R(t)=\int_{B_R} \big(u(t,x)-\bE[u(t,x)] \big)dx,
\]
where $B_R=\{x \in \bR^d;|x|<R\}$. We denote $\sigma_R^2(t)=\bE[F_R^2(t)]$. The goal of this article is to show that $F_R(t)/\sigma_R(t)$ converges in distribution to the standard normal distribution when $R \to \infty$, and give a precise rate of this convergence in the total variation distance. 
%A similar analysis has been done in the recent article \cite{BS26} for the stochastic heat %equation.

An effective approach to analyze the spatial average combines Malliavin calculus with Stein's method for normal approximation. This was developed in \cite{HNV20} for the one-dimensional nonlinear stochastic heat equation driven by space-time white noise, obtaining a quantitative central limit theorem and a functional central limit theorem. In \cite{HNVZ20} this was extended to the stochastic heat equation driven by noise that is white in time and has a Riesz spatial covariance. 

For the stochastic wave equation without drift, \cite{DNZ20} established quantitative and functional central limit theorems in dimension one for Gaussian noise that is white in time and fractional or white in space. In \cite{BNZ} the two-dimensional equation with a Riesz spatial covariance was treated, and \cite{NZ22} considered integrable spatial covariance kernels in dimensions one and two. A central ingredient in these works is a pointwise moment estimate for the first Malliavin derivative in terms of the fundamental solution of the wave equation. Establishing this estimate is particularly delicate in dimension two because the wave kernel is singular at the boundary of its support.

A complementary line of research concerns the hyperbolic Anderson model, for which $\sigma(u)=u$. In this case \cite{BNQZ} obtained moment estimates for increments of the solution's Malliavin derivatives, and applied a second-order Gaussian Poincar\'e inequality to prove quantitative and functional central limit theorems for spatially and temporally colored Gaussian noise. Rough spatial noises were subsequently considered in \cite{BY24} for the time-independent setting, and in \cite{BHWXY25} for noise that is white in time. Both of which also rely on moment estimates for the solution's Malliavin derivatives.

The works on the wave equation cited above do not include a nonlinear drift. Its presence producing a mixture of Lebesgue and stochastic integral terms which pose additional issues. For the one-dimensional stochastic heat equation driven by space-time white noise, \cite{BS26} recently established quantitative and functional central limit theorems in the presence of a nonlinear drift. The analysis uses a second-order Gaussian Poincar\'e inequality together with an estimate for products of heat kernels to control the second Malliavin derivative. The present article addresses the corresponding problem for the stochastic wave equation in dimensions one and two, allowing both integrable and Riesz spatial covariance kernels. In this regime, current methods for estimating the second Malliavin derivative are insufficient when the dimension is two. To overcome this, we instead estimate an integral of the second Malliavin derivative, which is compatible with the Gaussian Poincar\'e inequality. 
\medskip

To formulate our main results, we now specify the assumptions on the spatial covariance of the noise. We consider two cases: \\
(i) $f\in L^1(\bR^d)$ and (D) holds, or $d=1$ and the noise is white; if $d = 2$, we assume in addition that $f \in L^{\ell}(\bR^2)$ for some $\ell > 1$;
\\
(ii) $f(x)=|x|^{-\beta}$ is the Riesz kernel of order $\beta \in (0,d\wedge 2)$. In this case,
\[
\int_{\bR^d}\varphi(x)|x|^{-\beta}dx=c_{d,\beta} \int_{\bR^d}\cF \varphi(\xi) |\xi|^{-(d-\beta)}d\xi, \quad \mbox{for all} \quad \varphi \in \cS(\bR^d),
\]
where
\[
c_{d,\beta}=\pi^{-\frac{d}{2}} 2^{-\beta} \frac{\Gamma(\frac{d-\beta}{2})}{\Gamma(\frac{\beta}{2})}.
\]
The distinction between the integrable kernel and Riesz kernel is pervasive, as it determines many of the rates and normalizations throughout.

\medskip

Below are the main results of this article.


\begin{theorem}
\label{ergodic-th}
For any $t>0$, the process $\{u(t,x)\}_{x\in \bR^d}$ is ergodic.
\end{theorem}

\begin{theorem}[Limiting covariance]
\label{cov-th}
For any $t_1,t_2>0$,
\begin{equation}
\label{def-K}
K(t_1,t_2):=\lim_{R \to \infty}\frac{1}{R^{\gamma}}\bE[F_R(t_1) F_R(t_2)] \quad \mbox{is finite},
\end{equation}
where
\begin{equation}
\label{def-gamma}
\gamma:=\left\{
\begin{array}{ll} d & \mbox{in case (i)} \\
2d-\beta & \mbox{in case (ii) }.
\end{array} \right.
\end{equation}
In particular, $\sigma_R^2(t) \sim K(t,t) R^{\gamma}$ as $R \to \infty$, for any $t>0$.
\end{theorem}



\begin{theorem}[Quantitative Central Limit Theorem]
\label{QCLT}
For any $t>0$, there exists a constant $C_t>0$ depending on $t$ such that for any $R>0$,
\[
d_{TV}\left( \frac{F_R(t)}{\sigma_R(t)},Z\right) \leq C_t R^{\frac{\gamma-2d}{2}},
\]
% \[
% d_{TV}\left( \frac{F_R(t)}{\sigma_R(t)},Z\right) \leq 
% \left\{
% \begin{array}{ll} C_tR^{-d/2}, & \mbox{if $f \in L^1(\bR^d)$ and (D) holds, or $d=1$ and the noise is white } \\
% C_tR^{-\beta/2} & \mbox{if $f(x)=|x|^{-\beta}$ for some $\beta \in (0,d\wedge 2)$},
% \end{array} \right.
% \]
where $Z$ is a standard normal random variable, $\gamma$ is given by \eqref{def-gamma}, and $d_{TV}$ is the total variation distance.
\end{theorem}


\begin{theorem}[Functional Central Limit Theorem]
\label{FCLT}
%Let $\beta$ be given by \eqref{def-beta}.
The process $\{F_R(t)\}_{t\geq 0}$ has a $\delta$-H\"older continuous modification (denoted also $F_R$), for any $\delta \in (0,1/d)$. The process $\{R^{-\gamma/2} F_R(t)\}_{t\geq 0}$ converges in distribution in $C[0,\infty)$ as $R \to \infty$ to a zero-mean Gaussian process $\{\cG(t)\}_{t\geq 0}$ with covariance 
\[
\bE[\cG(t_1) \cG(t_2)]=K(t_1,t_2),
\]
where $\gamma$ is given by \eqref{def-gamma} and $K(t_1,t_2)$ is given by \eqref{def-K}.
\end{theorem}

The principal difficulty in extending the drift-free results to the present setting comes from the nonlinear drift. Although the pointwise estimate for the first Malliavin derivative remains valid, as shown in Theorem 3.1, the usual bound for the second derivative by a product of fundamental solutions cannot be used for a general nonlinear drift. Indeed, differentiating the drift twice produces a convolution involving $b''(u)D_{r,z}u\,D_{\theta,w}u$, which does not have the time-ordered product structure available in the linear Anderson setting. Instead of seeking an analogous pointwise bound, we exploit the compact support of $G$ and the resulting localization of the Malliavin derivatives to establish the spatially integrated moment estimate in Theorem 3.3. This estimate is sufficient, together with the first-derivative bound and the second-order Gaussian Poincar\'e inequality, to obtain the rate in Theorem \ref{QCLT}.

The drift also introduces difficulties in the limiting-covariance and functional-limit arguments. The covariance analysis requires separate treatment of the drift–drift and mixed contributions, while tightness requires temporal-increment estimates at the fluctuation scale $R^{\gamma/2}$. For the latter, we apply the Clark-Ocone formula to the centered drift increments, converting them into stochastic integrals whose moments can be controlled by the Burkholder–Davis–Gundy inequality and Theorem \ref{QCLT}. This yields bounds compatible with the normalization of the spatial average and completes the tightness argument.

\medskip

{\em Notation:} We use the notation $a_R \les b_R$ if $a_R \leq C b_R$ where $C$ is a constant that does not depend on $(a_R,b_R)$. We denote by $\| \cdot \|_p$ the norm in $L^p(\Omega)$ for any $p\geq 1$.
%% END Intro
%% BEGIN Prelims

\section{Preliminaries}

We first introduce the notation and auxiliary results that will be used throughout the paper.

\subsection{Inequalities}

The following is a collection of several Hardy--Littlewood--Sobolev--type inequalities that will be used throughout.

\medskip

If $f \in L^{\ell}(\bR^d)$ for some $\ell \geq 1$, using H\"older's inequality and Young's inequality, we have:
\begin{align}
\nonumber
\left| \int_{\bR^d}\int_{\bR^d}\varphi(x)\psi(y) f(x-y)dxdy \right| & \leq \|\varphi\|_{L^{\frac{2\ell}{2\ell-1}}(\bR^d)} \| \psi * f\|_{L^{2\ell}(\bR^d)} \\
\label{LHS-Lell}
& \leq \|\varphi\|_{L^{\frac{2\ell}{2\ell-1}}(\bR^d)} \|\psi\|_{L^{\frac{2\ell}{2\ell-1}}(\bR^d)} \|f\|_{L^{\ell}(\bR^d)}.
\end{align}



If $f(x)=|x|^{-\beta}$ with $\beta \in (0,d)$, %letting  $q=\frac{d}{2d-\beta}$,
by Hardy-Littlewood-Sobolev inequality, we have:
\begin{equation*}
\left|\int_{\bR^d} \int_{\bR^d}\varphi(x) \psi(y) |x-y|^{-\beta} dxdy\right| \leq c_{\beta,d} \|\varphi\|_{L^{\frac{2d}{2d-\beta}}(\bR^d)} \|\psi\|_{L^{\frac{2d}{2d-\beta}}(\bR^d)},
\end{equation*}
%for any $\varphi,\psi \in L^{2q}(\bR^d)$, 
where $c_{\beta,d}>0$ is a constant which depends on $(\beta,d)$; see e.g. Remark B.3 of \cite{BS26} for $d=1$, and Section 3.3 of \cite{BNZ} for $d=2$. 

To shorten the exposition, we combine the above inequalities in two distinct ways. First we have:
\begin{equation}
\label{LHS-q}
\left|\int_{\bR^d} \int_{\bR^d}\varphi(x) \psi(y) f(x-y) dxdy\right| \leq c_{f,d} \|\varphi\|_{L^{2q}(\bR^d)}\|\psi\|_{L^{2q}(\bR^d)},
\end{equation}
where $c_{f,d}>0$ is a constant that depends on $(f,d)$, and
\begin{equation}
\label{def-q}
q:=\left\{
\begin{array}{ll} 
1 & \mbox{if $d=1$ and $f \in L^1(\bR)$ or the noise is white}, \\
\frac{\ell}{2\ell-1} & \mbox{if $d = 2$ and $f \in L^1(\bR^2) \cap L^{\ell}(\bR^2)$ for some $\ell >1$}, \\
\frac{d}{2d-\beta} & \mbox{if $f(x)=|x|^{-\beta}$ for some $\beta \in (0,d)$}.
\end{array} \right.
\end{equation}
In the last two cases, $q \in (\frac12,1)$. Second, we have:
\begin{equation}
\label{LHS-q'}
\left|\int_{\bR^d} \int_{\bR^d}\varphi(x) \psi(y) f(x-y) dxdy\right| \leq c_{f,d} \|\varphi\|_{L^{2q'}(\bR^d)}\|\psi\|_{L^{2q'}(\bR^d)},
\end{equation}
where
\begin{equation}
\label{def-q'}
q':=\frac{d}{\gamma}=\left\{
\begin{array}{ll} 
1 & \mbox{if $f \in L^1(\bR^d)$ or $d=1$ and the noise is white}, \\
\frac{d}{2d-\beta} & \mbox{if $f(x)=|x|^{-\beta}$ for some $\beta \in (0,d)$}.
\end{array} \right.
\end{equation}
Inequalities \eqref{LHS-q} and \eqref{LHS-q'} are obviously the same, when $f(x)=|x|^{-\beta}$.
The subtle difference between them is that \eqref{LHS-q'} only requires $f \in L^1(\bR^d)$ even for $d=2$, i.e. we will use \eqref{LHS-q'} for arguments for which the hypothesis $f \in L^{\ell}(\bR^2)$ for some $\ell>1$ is {\em not} needed. Moreover, the exponent $q'$ on the right side of \eqref{def-q'} depends on the parameter $\gamma$, which appears in the QCLT rate, and in the normalization in the FCLT. 

\subsection{Properties of $G$}

In this section, we include some properties of the wave kernel $G$. First, note that
\begin{equation}
\label{int G p}
 \int_{\bR^d} G^p_t(x) dx = 
 \begin{cases} 
  2^{1-p}t & \mbox{for any $p >0$, if $d=1$} \\
  \frac{(2\pi)^{1-p}}{2-p}t^{2-p} & \mbox{for any $p\in (0, 2)$, if $d=2$}. \\
 \end{cases}
\end{equation}
The next result gives an explicit expression for the convolution of two wave kernels.
\begin{lemma}
\label{lemG*G}
Let $d\in \{1,2\}$. For any $t,s>0$ and $x\in \bR^d$, we have:
\begin{equation}
\label{G*G}
(G_t*G_s)(x)=\frac{1}{2}\int_{|t-s|}^{t+s} G_{\theta}(x) d\theta.
\end{equation}
\end{lemma}

\begin{proof}
We show that both sides have the same Fourier transforms. On the left side,
\begin{align*}
\cF(G_t*G_s)(\xi)&= \cF G_t(\xi)\cF G_s(\xi)=\frac{\sin(t|\xi|) \sin(s|\xi|)}{|\xi|^2}.
\end{align*}
On the right side,
\begin{align*}
\frac{1}{2}\int_{|t-s|}^{t+s} \cF G_{\theta}(\xi) d\theta=\frac{1}{2}\int_{|t-s|}^{t+s} \frac{\sin(\theta|\xi|)}{|\xi|} d\theta=\frac{\cos\big(|t-s||\xi|\big)- \cos \big( (t+s)|\xi|\big)}{2|\xi|^2}.
\end{align*}
\end{proof}
For any $t>0$ and $x \in \bR^d$, we denote $K_t(x) := (G_t*G_t)(x)$. By Lemma \ref{lemG*G},
\[
K_t(x) =\frac{1}{2}\int_0^{2t} G_{\theta}(x)d\theta. 
\]
Hence, $t \mapsto K_t(x)$ is non-decreasing on $\bR_{+}$, and $K_t(x)=0$ if $|x| \geq 2t$. By direct calculation,
\begin{align}
\label{def-K1}
 K_t(x) & = \frac{1}{4} (2t - |x|) \, 1_{\{|x| < 2t\}}  \qquad \text{if} \ d = 1 \\
 \label{def-K2}
 K_t(x) & = \frac{1}{4\pi}  \log\left(\frac{2t+\sqrt{4t^2-|x|^2}}{|x|}\right)  1_{\{|x|<2t\}}\qquad \text{if} \ d = 2.
\end{align}
The following lemma provides an upper bound similar to inequality \eqref{LHS-q}, but for the kernel $fK_t^{\alpha}$ replacing $f$.
\begin{lemma} \label{fK int}
Let $ t \in [0, T]$ and $\alpha>0$. Then 
\begin{enumerate}
 \item[(i)] If $d = 1$ and $f \in L^1(\bR)$ or the noise is white, then
\[
\int_{\bR} \int_{\bR} \varphi(x)\psi (y) f(x-y) K_t^{\alpha}(x-y) dx dy  \leq T^{\alpha} \|f\|_{L^1(\bR)} \|\varphi\|_{L^{2}(\bR)}\|\psi\|_{L^{2}(\bR)}.
\]
\item[(ii)] If $d\in \{1,2\}$ and either (a) $f \in L^\ell(\bR^d)$ for some $\ell >1$, or (b) $f(x) = |x|^{-\beta}$ for some $\beta \in (0,d)$, then 
\[
\int_{\bR^d} \int_{\bR^d} \varphi(x)\psi(y) f(x-y)K_t^{\alpha}(x-y)  dx dy \leq C_{T,\alpha,p',q} \|\varphi\|_{L^{2p'}(\bR^d)}\|\psi\|_{L^{2p'}(\bR^d)},
\]
for any $p'\in (q,1)$, where $q=\frac{\ell}{2\ell-1}$ in case (a), and $q=\frac{d}{2d-\beta}$ in case (b). 
\end{enumerate}
\end{lemma}

\begin{proof} 
(i) By \eqref{def-K1}, it follows that $K_t(x) \leq T$ for any $x \in \bR$ and $t \in [0,T]$. Therefore, 
\[
 \int_{\bR} \int_{\bR} \varphi(x) \psi(y)  f(x-y)K_t^{\alpha}(x-y) dx dy \leq  T^{\alpha} \int_{\bR} \int_{\bR} \varphi(x) \psi(y) f(x-y) dx dy.
\]
The conclusion follows by applying \eqref{LHS-Lell} with $\ell=1$.  

(ii) Using \eqref{def-K1} and \eqref{def-K2}, it follows that for all $\varepsilon > 0$,
\begin{equation}
\label{bound-K}
 K_t(x) \leq C_{T,\epsilon} |x|^{-\varepsilon}.
\end{equation}

We fix $p' \in (q,1)$. We apply H\"older's inequality, followed by Young's inequality:
\begin{align*}
I &:=\int_{\bR^d} \int_{\bR^d}  \varphi(x) \psi(y) f(x-y)K_t^{\alpha}(x-y) dx dy 
= \int_{\bR^d} (\varphi * fK_t^{\alpha})(y) \psi(y) dy \\
& \quad \leq \| \varphi* fK_t^n  \|_{L^a(\bR^d)} \|\psi \|_{L^{2p'}(\bR^d)}  \quad \text{with}\  \frac{1}{a} + \frac{1}{2p'} = 1, \\
& \quad \leq  \| fK_t^n\|_{L^{b}(\bR^d)}  \|\varphi \|_{L^{2p'}(\bR^d)}\|\psi \|_{L^{2p'}(\bR^d)} \quad  \text{with } \  \frac{1}{b} +\frac{1}{2p'} = 1 + \frac{1}{a} = 2 - \frac{1}{2p'} .
\end{align*}
Note that $b=\frac{p'}{2p'-1}>1$ since $p'\in (\frac12,1)$. 

To estimate $\|K_t^{\alpha}\|_{L^{c}(\bR^d)}$, we consider separately the two cases.
In case (a), we apply H\"older's inequality with $\frac{1}{p_1}+\frac{1}{p_2}=1$ and $p_1=\frac{\ell}{b}>1$. (This means that $\frac{1}{\ell}<\frac{1}{b}=2 - \frac{1}{p'}$, i.e. $p'>\frac{\ell}{2\ell -1}=q$.) Then
\begin{align*}
\| fK_t^{\alpha}\|_{L^{b}(\bR^d)}  & \leq \| f\|_{L^\ell(\bR^d)} \|K_t^{\alpha}\|_{L^{c}(\bR^d)} \qquad \text{with } \frac{1}{l} + \frac{1}{c} = \frac{1}{b}.
\end{align*}
We use \eqref{bound-K} and the compact support of $K_t$. Then,
$\|K_t^{\alpha}\|_{L^{c}(\bR^d)} \leq C_{T,\varepsilon}\int_{|x| < 2t} |x|^{-\alpha \varepsilon c} dx < \infty$, provided that we choose $\varepsilon > 0$ so that $\alpha \varepsilon  c< d$. This implies that
\[
I \leq C_{T,\varepsilon,\alpha,c,d} \| f\|_{L^\ell(\bR^d)} \|\varphi\|_{L^{2p'}(\bR^d)}\|\psi\|_{L^{2p'}(\bR^d)}. 
\]

In case (b), using \eqref{bound-K} and the compact support of $K_t$, we have:
\[
\|fK_t^{\alpha}\|_{L^b}^b \leq C_{T,\varepsilon} \int_{\{|x|<2t\}}  |x|^{-\beta b-\alpha\varepsilon  b} dx<  \infty,
\]
provided that we choose $\varepsilon>0$ such that $(\beta+\alpha\varepsilon ) b< d$. This means that 
$$2-\frac{1}{p'}=\frac{1}{b}>\frac{\beta+\alpha \varepsilon}{d} \quad \mbox{i.e.} \quad \frac{\alpha \varepsilon}{d}< \frac{2d-\beta}{d}-\frac{1}{p'}.$$
This is possible since $p'>\frac{d}{2d-\beta} =q$. Then,
\[
I \leq C_{T,\varepsilon,\alpha,\beta,b,d} \|\varphi\|_{L^{2p'}(\bR^2)}\|\psi\|_{L^{2p'}(\bR^2)}. 
\]
\end{proof}

\bigskip 

We also have the follwing inequalities for the wave kernel $G$ integrated over a ball of radius $R$. 
For any $R>0$, $0\leq r \leq t$ and $y \in \bR^d$, we denote:
\begin{equation}
\label{def-var}
\varphi_{t,R}(r,y)=\int_{B_R} G_{t-r}(x-y)dx.
\end{equation}

Observe that, $0 \leq \varphi_{t,R}(r,y) \leq t-r$ and $\varphi_{t,R}(r,y)$ is supported on $\{|y| \leq R + t - r \}$. In particular, when $R \geq t$, $\varphi_{t,R}(r,y)$ is supported on $B_{2R}$. Therefore, for any $p\geq 1$, $0\leq r\leq t$ and $R>t$,
\begin{equation}
\label{est}
 \| \varphi_{t,R}(r,\cdot) \|_{L^p(\bR^d)} \les R^{d/p} (t-r).
\end{equation}
From the proof of Lemma 4.3 in \cite{NZ22}, we know that
\[
\Big|\varphi_{t,R}(r,y) - \varphi_{s,R}(r,y) \Big| \les (t-s)^{1/d} \mathbf 1_{\{|y|\leq R+t\}} \quad \mbox{for any $0\leq r\leq s\leq t$}.
\]
As a consequence, for any $p\geq 1$, $0\leq r\leq s\leq t$ and $R>t$,
\begin{equation}
\label{diff est}
 \| \varphi_{t,R}(r,\cdot) - \varphi_{s,R}(r,\cdot)\|_{L^p(\bR^d)} \les R^{d/p} (t-s)^{1/d}.
\end{equation}

\subsection{Malliavin Calculus}

In this section, we include some background material about Malliavin calculus.

\medskip

Let $\bL^{1,2}$ be the class of processes $u \in L^2(\Omega;\cH)$ such that $u(t,x)\in \bD^{1,2}$ for all $(t,x)\in \bR_{+}\times \bR^d$, and there exists a measurable modification of
$\{D_{t,x}u(s,y);(t,x),(s,y)\}$ such that
\[
\bE \int_{(\bR_+ \times \bR^d)^2} D_{r,y}u(t,x)D_{r,y'}u(t,x') f(x-x')f(y-y') dxdx'dydy'drds < \infty.
\]

The operators $D$ and $\delta$ satisfy the following Heisenberg commutation rule.
\begin{proposition}[Proposition 1.3.8 of \cite{nualart06}] If $u \in \bL^{1,2}$, $\{D_{t,x}u(s,y)\}_{(s,y)} \in {\rm Dom}(\delta)$ for almost all $(t,x)$, and there exists a measurable modification of $\{\int D_{t,x}u(s,y)W(\delta s,\delta y)\}_{(t,x)}$ which is in $L^2(\Omega;\cH)$, then $\delta(u) \in \bD^{1,2}$ and the following identity holds, 
\begin{equation}
\label{Heisenberg}
D(\delta(u))=u+\delta(Du).
\end{equation}
\end{proposition}

The Malliavin derivative satisfies the following chain rule.

\begin{proposition}[Proposition 1.2.3 of \cite{nualart06}] 
For any $F \in \bD^{1,2}$ and $\varphi \in C^1(\bR)$ with $\varphi'$ bounded,
\begin{equation}
\label{chain}
D\varphi(F)=\varphi'(F)DF.
\end{equation}
\end{proposition}

Similarly to Corollary 1.2.1 of \cite{nualart06}, if $F\in \bD^{1,2}$ is $\cF_t$-measurable, then 
\begin{equation}
\label{Mal-zero}
D_{r,z}F=0 \quad \mbox{for all $r>t$ and $z \in \bR^d$}.
\end{equation}

Finally, we have the following version of Poincar\'e inequality (see e.g. relation (2.3) of \cite{NZ22}).

\begin{proposition}
For any $F,G \in \bD^{1,2}$,
\begin{equation}
\label{Poincare}
|{\rm Cov}(F,G)| \leq \int_{\bR_+} \int_{\bR^{2d}} \|D_{r,y}F \|_2 \|D_{r,y'}G\|_2 f(y-y') dydy'dr.
\end{equation}
\end{proposition}


%% END Prelims
%% BEGIN Key Estimates
\section{Estimates on Malliavin Derivatives}

In this section, we provide some key estimates for the Mallivain derivatives of the solution $u$, which will be used in the proofs of the main results.

We denote by $(u_n)_{n\geq 0}$ the sequence of Picard iterations, given by $u_0(t,x)=1$ and
\begin{equation}
\label{def-Picard}
u_{n+1}(t,x)=1+\int_0^t \int_{\bR^d}G_{t-s}(x-y)b(u_n(s,y))dyds+\int_0^t \int_{\bR^d}G_{t-s}(x-y)\s(u_n(s,y))W(ds,dy).
\end{equation}


\begin{theorem}
\label{key-Du}
For any $0\leq r\leq t \leq T$, $x,z \in \bR^d$ and $p\geq 2$,
\[
\|D_{r,z}u(t,x)\|_p \leq C_{T,p} G_{t-r}(x-z),
\]
where $C_{T,p}>0$ is a constant that depends on $(T,p,q,\|b'\|_\infty,\|\sigma'\|_\infty)$.
\end{theorem}


\begin{proof}
Using Heisenberg commutation rule \eqref{Heisenberg} followed by the chain rule \eqref{chain}, we infer that the sequence $(Du_n)_{n\geq 0}$ satisfies the following recurrence relation:
\begin{align}
\nonumber
D_{r,z}u_{n+1}(t,x)&=G_{t-r}(x-z)\s\big(u_n(r,z)\big)+\int_r^t \int_{\bR^d}G_{t-s}(x-y) b'\big(u_n(s,y)\big)D_{r,z}u_n(s,y) dyds+\\
\label{rec-Dun}
&\quad \int_r^t \int_{\bR^d} G_{t-s}(x-y)\s'\big(u_n(s,y)\big) D_{r,z}u_n(s,y) W(ds,dy),
\end{align}
for any $n\geq 0$. Note that $D_{r,z}u_{0}(t,x)=0$ since $u_0(t,x)=1$. 

We use Burkholder-Davis-Gundy (BDG) inequality and Minkowski inequality to estimate the $p$-th moment of the stochastic integral. Using \eqref{def-K}, it follows that for any $0<r<t\leq T$, $x,z \in \bR^d$ and $n\geq 0$,
\begin{align}
\nonumber
&\|D_{r,z}u_{n+1}(t,x)\|_p^2  \leq B_{t} \left\{ G_{t-r}^2(x-z)+  \left(\int_r^t \int_{\bR^d} G_{t-s}(x-y)\|D_{r,z}u_n(s,y)\|_pdyds\right)^2 \right.\\	
\nonumber
& \quad +\left.\int_r^t \int_{(\bR^d)^2} G_{t-s}(x-y) G_{t-s}(x-y') \|D_{r,z}u_n(s,y)\|_p \|D_{r,z}u_n(s,y')\|_p  f(y-y') dydy'ds \right\} \\\
\label{recD1}
& \quad \quad \quad \quad \quad \quad  =: B_{t} \{ G_{t-r}^2(x-z)+  I+II \},
\end{align}
where $B_{t}$ is a constant that depends on $(t,p,\|b'\|_{\infty},\|\sigma'\|_{\infty})$ and is non-decreasing in $t$.

We will prove below that for any $0<r<t\leq T$, $x,z \in \bR^d$ and $n\geq 1$,
\begin{equation}
\label{key-un}
\|D_{r,z}u_{n}(t,x)\|_p \leq C_{T,p} G_{t-r}(x-z).
\end{equation}

The conclusion will follow by Lemma A.1 of \cite{BS26}, since $\{Du_n(t,x)\}_{n\geq 1}$ converges to $Du(t,x)$ in the weak topology $L^2(\Omega;\cH)$ (by Lemma 1.2.3 of \cite{nualart06}).

\medskip

{\em Case 1.} $d=1$.
For the first integral on the RHS of \eqref{recD1}, we use H\"older's inequality with respect to the measure $\mu_t(ds,dy)=G_{t-s}(x-y)dsdy$ on $[r,t] \times \bR$ whose total mass is $(t-r)^2/2$. For the second integral of \eqref{recD1}, we use \eqref{LHS-q}. We obtain that for any $0<r<t\leq T$, $x,z \in \bR, n\geq 0$,
\begin{align*}
\|D_{r,z}u_{n+1}(t,x)\|_p^2  & \leq  B_{T} \left\{ G_{t-r}^2(x-z)+ \frac{T^2}{2}\int_r^t \int_{\bR} G_{t-s}^2(x-y) \|D_{r,z}u_n(s,y)\|_p^2 dyds \right.\\
& \quad \quad \quad \left. c_{f,1}\int_r^t \left( \int_{\bR}G_{t-s}^{2q}(x-y) \|D_{r,z}u_n(s,y)\|_p^{2q}dy\right)^{1/q}ds \right\},
\end{align*}
where $q$ is given by \eqref{def-q}. Using the fact that $G_t^{2q}(x)=c_{q} G_t^{2}(x)$ and applying H\"older's inequality with respect to the measure $G_{t-s}^2(x-y)dy$ whose total mass is $(t-s)/2$, it follows that:
\[
\left(\int_{\bR} G_{t-s}^{2q}(x-y) \|D_{r,z}u_n(s,y)\|_p^{2q} dy\right)^{1/q} \leq c_{q}^{q} \frac{T}{2}\int_{\bR} G_{t-s}^2(x-y)\|D_{r,z}u_n(s,y)\|_p^{2} dy.
\]
Hence
\begin{align*}
&\|D_{r,z}u_{n+1}(t,x)\|_p^2  \leq  A_T \left\{ G_{t-r}^2(x-z)+ \int_r^t \int_{\bR} G_{t-s}^2(x-y) \|D_{r,z}u_n(s,y)\|_p^2 dyds \right\},
\end{align*}
where $A_T=B_{T} \big(1+\frac{T^2}{2}+c_{q}^{q} \frac{T}{2}\big)$.
By induction on $n\geq 1$, it follows that 
\[
\|D_{r,z}u_{n}(t,x)\|_p^2 \leq A_T \left\{ G_{t-r}^2(x-z)+ \sum_{k=1}^{n-1} A_T^k I_k(t,x,r,z)\right\},
\]
for any $0<r<t\leq T$ and $x,z \in \bR$, where
\begin{align}
\label{def-Ik}
I_k(t,x,r,z)&=\int_{r<t_1<\ldots<t_k<t} \int_{\bR^k}\prod_{i=1}^{k}G_{t_{i+1}-t_i}^2(x_{i+1}-x_i) G_{t_1-r}^2(x_1-z)d\pmb{x}_k d \pmb{t}_k,
\end{align}
with $\pmb{x}_k=(x_1,\ldots,x_k)$, $\pmb{t}_k=(t_1,\ldots,t_k)$, $t_{k+1}=t$ and $x_{k+1}=x$. Using the fact that $G_t^2(x)=\frac{1}{2}G_t(x)$ and inequality
\[
\int_{\bR}G_{t-s}(x-y)G_{s-r}(y-z)dy \leq (t-r)G_{t-r}(x-z),
\]
it follows that 
\[
I_k(t,x,r,z)\leq \frac{(t-r)^{2k}}{2^k k!}G_{t-r}(x-z).
\]
Hence, for any $0<r<t\leq T$, $x,z \in \bR$ and $n\geq 1$,
\[
\|D_{r,z}u_{n}(t,x)\|_p^2 \leq 2A_T \left(\sum_{k\geq 0} A_T^k \frac{T^{2k}}{2^k k!} \right) G_{t-r}^2(x-z)=:C_{T,p}^2 G_{t-r}^2(x-z).
\]

\medskip

{\em Case 2.} $d=2$. We proceed similarly to Lemma A.2 of \cite{BS26}. We fix $r \in [0,T]$ and $z \in \bR^2$.
By induction on $n \geq 1$,  $D_{r,z}u_n(t,x)=0$ if $|x-z|\geq t-r$. We denote 
\[
f_n(t,x)=\|D_{r,z}u_n(t,x)\|_p \quad \mbox{and} \quad \widetilde{f}_n(t)=\sup_{|x-z|<t-r}\frac{f_n(t,x)}{G_{t-r}(x-z)}.
\]

We suppose first that $\widetilde{f}_n(t)<\infty$ for all $t \in [r,T]$, and we derive a recurrence relation for $\widetilde{f}_{n+1}$. 

We bound the terms on the right-hand side of \eqref{recD1}. We start with $I$, the term due to the drift. We have:
\begin{align*}
\int_r^t \int_{\bR^2} G_{t-s}(x-y) f_n(s,y)dyds & \leq \int_r^t \widetilde{f}_n(s) \int_{\bR^2} G_{t-s}(x-y) G_{s-r}(y-z) dy ds \\
& = \frac{1}{2}\int_r^t  \widetilde{f}_n(s) \int_{|t-2s+r|}^{t-r} G_{\theta}(x-z)d\theta ds,
\end{align*}
where for the second line we use the  identity \eqref{G*G}.


Considering separately the cases $\theta \in [r,\frac{t+r}{2}]$ and $\theta \in [\frac{t+r}{2},t-r]$, and using Fubini's theorem, it follows that
\[
\int_r^t \int_{\bR^2} G_{t-s}(x-y) f_n(s,y)dyds \leq \frac{1}{2} \int_0^{t-r} \left(\int_{\frac{t+r-\theta}{2}}^{\frac{t+r+\theta}{2}} \widetilde{f}_n(s) ds \right) G_{\theta}(x-z)d\theta.
\]
Using Caucy-Schwarz inequality, we obtain that for any $\theta \in [0,t-r]$,
\begin{align*}
\int_{\frac{t+r-\theta}{2}}^{\frac{t+r+\theta}{2}} \widetilde{f}_n(s) ds  & \leq \theta^{1/2}
\left(\int_{\frac{t+r-\theta}{2}}^{\frac{t+r+\theta}{2}} \widetilde{f}_n^2(s) ds\right)^{1/2} \leq
\theta^{1/2}
\left(\int_{r}^{t} \widetilde{f}_n^2(s) ds\right)^{1/2}.
\end{align*}
Hence,
\[
\int_r^t \int_{\bR^2} G_{t-s}(x-y) f_n(s,y)dyds \leq \frac{1}{2} \left(\int_{r}^{t} \widetilde{f}_n^2(s) ds\right)^{1/2}\int_0^{t-r} \theta^{1/2} G_{\theta}(x-z)d\theta.
\]
By direct calculation,
\begin{equation}
\label{theta-G}
\int_0^t \theta^{1/2} G_{\theta}(x) d\theta \leq 2t^{3/2} G_{t}(x) \quad \mbox{for any $t>0$}.
\end{equation}
Indeed,
\begin{align*}
\int_0^t \theta^{1/2} G_{\theta}(x) d\theta &=\frac{1}{2\pi}\int_{|x|}^t \sqrt{\frac{\theta}{\theta+|x|}} \cdot \frac{1}{\sqrt{\theta-|x|}} d\theta \leq \frac{1}{2\pi} \int_{0}^{t-|x|} \frac{1}{\sqrt{\eta}} d\eta=\frac{1}{\pi}\sqrt{t-|x|} \\
& \leq 2 \sqrt{(t-|x|)(t^2-|x|^2)} G_t(x) \leq 2t^{3/2} G_{t}(x).
\end{align*}
It follows that
\begin{equation}
\label{bound-I}
I \leq   (t-r)^{3} G_{t-r}^2(x-z) \int_{r}^{t} \widetilde{f}_n^2(s) ds.
\end{equation}

We now bound the term $II$, due to the noise. Using the definition of $\widetilde{f}_n(t)$ and inequality \eqref{LHS-q} with $q$ given by \eqref{def-q}, we have:
\begin{align*}
II & \leq \int_r^t \widetilde{f}_n^2(s) \int_{\bR^2} \int_{\bR^2} G_{t-s}(x-y) G_{t-s}(x-y') G_{s-r}(y-z) G_{s-r}(y'-z) f(y-y') dydy'ds \\
& \leq c_{f,2} \int_r^t \widetilde{f}_n^2(s) \big[ G_{t-s}^{2q}* G_{s-r}^{2q}(x-z)\big]^{1/q}ds.  
\end{align*}
We apply H\"older's inequality, with $\frac{1}{p}+\frac{1}{p'}=1$:
\begin{align*}
II \le c_{f,2} \left(\int_r^t \widetilde{f}_n^{2p}(s)ds \right)^{1/p} \left( \int_r^t \big[ G_{t-s}^{2q}* G_{s-r}^{2q}(x-z)\big]^{p'/q}ds \right)^{1/p'}.
\end{align*}
We now apply Lemma 4.3 of \cite{BNZ} with $\delta=p'/q$, which means that we must choose $1<p'<\frac{q}{2q-1}$. Note that this lemma remains valid for any $\delta(1-2q)+1>0$ and $q \in (\frac{1}{2},1)$. We obtain that
\begin{equation}
\label{bound-II}
II \leq C_{f,q}^2  (t-r)^{1-\delta(1-2q)} G_{t-r}^{\delta(2q-1)}(x-z) \left(\int_r^t \widetilde{f}_n^{2p}(s)ds \right)^{1/p},
\end{equation}
for some constant $C_{f,q}>1$.

We put together the bounds from \eqref{bound-I} and \eqref{bound-II}. First, by H\"older's inequality, 
\[
\int_r^t \widetilde{f}_n^2(s) ds \leq (t-r)^{1/p'} \left(\int_r^t \widetilde{f}_n^{2p}(s)ds \right)^{1/p}.
\]
Secondly, we use the fact that $G_{t-r}^{\delta(2q-1)}(x-z) \leq (t-r)^{2-\delta(2q-1)} G_{t-r}^{2}(x-z)$, due to the trivial bound:
\begin{equation}
\label{G-ab}
G_t^a(x) \leq t^{b-a} G_t^b(x) \quad \mbox{for any $t>0$ and $a<b$}.
\end{equation}
Hence,
\[
f_{n+1}^2(t,x) \le C_{f,q}^2 G_{t-r}^2(x-z) \left\{1+ [(t-r)^{3+\frac{1}{p'}}+(t-r)^{3-2\delta(2q-1)}] \left( \int_r^t \widetilde f_n^{2p}(s)ds\right)^{1/p}\right\}
\]
We divide by $G_{t-r}^2(x-z)$ and we take the supremum over $x$ with $|x-z|<t-r$. Then, we take power $p$. We obtain that:
\[
\widetilde{f}_{n+1}^{2p}(t) \le C_{f,q}^{2p} \left\{1+ [(t-r)^{3+\frac{1}{p'}}+(t-r)^{3-2\delta(2q-1)}]  \int_r^t \widetilde f_n^{2p}(s)ds\right\}.
\]

This relation is similar to (A.10) of \cite{BS26}. Proceeding as in Step 2 and 3 of the proof of Lemma A.1 of \cite{BS26}, we conclude that $\widetilde{f}_{n}(t) \leq C_{f,q} C_t$ where $C_t$ is a constant that is non-decreasing in $t$.
\end{proof}

\begin{corollary}
\label{D2u-support}
Let $0\leq t \leq T$, $r,\theta \in [0,t]$, $x,z,w \in \bR^d$, and $p \geq 2$. If $|x-z| > t-r$ or $|x-w| > t - \theta$. Then, 
\[
\|D_{(r,z),(\theta,w)}^2u(t,x)\|_p = 0. 
\]
\end{corollary}

Recall definition \eqref{def-q} of $q$.
%Let $p'=1$ be such that,
%\begin{equation}
%\label{def-p'}
% p'= 
% \begin{cases}
% 1 & \text{when } q = 1 \\
% q < p' < 1 & \text{otherwise.}
% \end{cases}
%\end{equation}


\begin{theorem}
\label{key-D2u-int}
For any $0\leq r,\theta \leq t' \leq T$, $w \in \bR^d$,  and $p\geq 2$,
\[
\int_{\bR^d} \left[ \sup_{r \vee \theta \leq t \leq t'} \int_{\bR^d} \|D_{(r,z),(\theta,w)}^2u(t,x)\|^{2p'}_p dx \right] dz \leq C_{T,p,p'},
\]
where $p'=1$ if $d=1$ and $f \in L^1(\bR)$ or the noise is white, and  $p' \in (q,1)$ is arbitrary if $d=2$ and $f \in L^1(\bR^2) \cap L^{\ell}(\bR^2)$ for some $\ell>1$, or $f(x)=|x|^{-\beta}$ for some $\beta \in (0,d)$.

\end{theorem}
\begin{proof}
We first note that $D_{r,z}u(t,x) = 0$ whenever $t < r$ by \eqref{Mal-zero}. Furthermore, we use the convention $G_t(x) = 0$ whenever $t < 0$. Let $0\leq r,\theta \leq t \leq t' \leq T$. We denote $\rho = r \vee \theta$.

Taking the first and then second Malliavin derivative of the mild solution \eqref{mild soln} using \eqref{chain} and \eqref{Heisenberg}, 
\begin{align*}
D^2_{(r,z),(\theta,w)}u(t,x) & =  G_{t-r}(x-z) \s'\big(u(r,z)\big) D_{\theta,w} u(r,z) \\
& + G_{t-\theta}(x-w) \s'\big(u(\theta,w)\big) D_{r,z} u(\theta,w) \\
& + \int_\rho^t \int_{\bR^d} G_{t-s}(x-y) b''\big(u(s,y)\big) D_{r,z}u(s,y) D_{\theta,w}u(s,y)  dyds \\
& + \int_\rho^t \int_{\bR^d} G_{t-s}(x-y) b'\big(u(s,y)\big) D_{(r,z),(\theta,w)}^2 u(s,y) dyds \\
& + \int_\rho^t \int_{\bR^d} G_{t-s}(x-y) \s''\big(u(s,y)\big) D_{r,z}u(s,y) D_{\theta,w}u(s,y)  W(ds,dy)\\
& + \int_\rho^t \int_{\bR^d} G_{t-s}(x-y) \s'\big(u(s,y)\big) D_{(r,z),(\theta,w)}^2 u(s,y) W(ds,dy).
\end{align*}
Using BDG inequality and Minkowski's inequality,
\begin{align*}
\|D^2_{(r,z),(\theta,w)}u(t,x)\|_p^{2p'} & \les  G^{2p'}_{t-r}(x-z) \|\s'\big(u(r,z)\big) D_{\theta,w} u(r,z)\|_p^{2p'} \\
& + G^{2p'}_{t-\theta}(x-w) \|\s'\big(u(\theta,w)\big) D_{r,z} u(\theta,w)\|_p^{2p'} \\
& + \left(\int_\rho^t \int_{\bR^d} G_{t-s}(x-y) \|b''\big(u(s,y)\big) D_{r,z}u(s,y) D_{\theta,w}u(s,y)\|_p  dyds\right)^{2p'} \\
& + \left(\int_\rho^t \int_{\bR^d} G_{t-s}(x-y) \|b'\big(u(s,y)\big) D_{(r,z),(\theta,w)}^2 u(s,y)\|_p dyds\right)^{2p'} \\
& + \bigg(\int_\rho^t \int_{\bR^d} \int_{\bR^d} G_{t-s}(x-y) \|\s''\big(u(s,y)\big) D_{r,z}u(s,y) D_{\theta,w}u(s,y)\|_p f(y-y')  \\
& \qquad G_{t-s}(x-y') \|\s''\big(u(s,y')\big) D_{r,z}u(s,y') D_{\theta,w}u(s,y')\|_p dy dy' ds \bigg)^{p'}\\
& + \bigg( \int_\rho^t \int_{\bR^d} \int_{\bR^d} G_{t-s}(x-y) \|\s'\big(u(s,y)\big) D_{(r,z),(\theta,w)}^2 u(s,y)\|_pf(y-y')  \\
& \qquad G_{t-s}(x-y') \|\s'\big(u(s,y')\big) D_{(r,z),(\theta,w)}^2 u(s,y')\|_p dy dy' ds \bigg)^{p'}.
\end{align*}
Using Theorem \ref{key-Du}, H\"older's inequality, and the boundedness of $b'$,$b''$,$\s'$ and $\s''$, 
\[
 \int_{\bR^d} \sup_{\rho \leq t \leq t'} \int_{\bR^d} \|D^2_{(r,z),(\theta,w)}u(t,x)\|_{p}^{2p'} dx dz \les \sum_{i=1}^5 \cB_i
\]
where
\begin{align*}
 \cB_1 & = \int_{\bR^d} \sup_{\rho \leq t \leq t'} \int_{\bR^d} \Big(G^{2p'}_{t-r}(x-z)G^{2p'}_{r-\theta}(z-w) + G^{2p'}_{t-\theta}(x-w)G^{2p'}_{\theta-r}(w-z) \Big) dx dz, \\
 \cB_2 & = \int_{\bR^d} \sup_{\rho \leq t \leq t'} \int_{\bR^d} \left( \int_\rho^t \int_{\bR^d} G_{t-s}(x-y)G_{s-r}(y-z)G_{s-\theta}(y-w) dy ds\right)^{2p'} dx dz, \\
 \cB_3 & = \int_{\bR^d} \sup_{\rho \leq t \leq t'}  \int_{\bR^d} \left( \int_\rho^t \int_{\bR^d} G_{t-s}(x-y) \|D^2_{(r,z),(\theta,w)}u(s,y)\|_{p}  dy ds\right)^{2p'} dx dz, \\
 \cB_4 & = \int_{\bR^d} \sup_{\rho \leq t \leq t'}  \int_{\bR^d} \bigg( \int_\rho^t \int_{\bR^d} \int_{\bR^d} G_{t-s}(x-y)G_{s-r}(y-z)G_{s-\theta}(y-w) f(y-y') \\
 & \quad G_{t-s}(x-y')G_{s-r}(y'-z)G_{s-\theta}(y'-w) dydy' ds\bigg)^{p'} dx dz, \\
 \cB_5 & = \int_{\bR^d} \sup_{\rho \leq t \leq t'}  \int_{\bR^d} \bigg( \int_\rho^t \int_{\bR^d} \int_{\bR^d} G_{t-s}(x-y) \|D^2_{(r,z),(\theta,w)}u(s,y)\|_{p} f(y-y') \\
 & \quad G_{t-s}(x-y') \|D^2_{(r,z),(\theta,w)}u(s,y')\|_{p} dy dy' ds\bigg)^{p'} dx dz. 
\end{align*}

We compute each term individually. For $\cB_1$ we use \eqref{int G p}, noting that $2p'<2$ if $d=2$:
% which is then bounded by powers of $T$. Therefore,
\[
\cB_1 \les \int_{\bR^d} \Big( G^{2p'}_{r-\theta}(z-w) + G^{2p'}_{\theta-r}(w-z)\Big) dz  \les C_T.
\]

For $\cB_2$ we use H\"older's inequality with respect to the finite measure $\mathbf 1_{\{\rho < s < t\}} \mathbf 1_{\{|x-y| \leq T\}}dyds$ and exponent $2p' > 1$:
\[
 \cB_2 \les \int_{\bR^d} \sup_{\rho \leq t \leq t'}  \int_\rho^t \int_{\bR^d} \int_{\bR^d}  G^{2p'}_{t-s}(x-y)G^{2p'}_{s-r}(y-z)G^{2p'}_{s-\theta}(y-w) dx dy ds dz. \\
\]
Applying \eqref{int G p} with respect to $dx$, 
\[
 \cB_2 \les \int_{\bR^d} \sup_{\rho \leq t \leq t'}  \int_\rho^t \int_{\bR^d}G^{2p'}_{s-r}(y-z)G^{2p'}_{s-\theta}(y-w) dy ds dz \leq \int_\rho^{t'} \int_{\bR^d} \int_{\bR^d} G^{2p'}_{s-r}(y-z)G^{2p'}_{s-\theta}(y-w) dz dy ds.
\]
 Applying \eqref{int G p} with respect to $dz$ and $dy$ then integrating in $ds$, 
\[
 \cB_2 \les \int_\rho^T \int_{\bR^d} \int_{\bR^d} G^{2p'}_{s-\theta}(y-w) dy ds \les \int_\rho^T ds \les C_T.
\]

Similarly for $\cB_3$, we use H\"older's inequality with respect to the measure $\mathbf 1_{\{\rho \leq s \leq t\}} \mathbf 1_{\{|x-y| \leq T\}}dyds$:
\[
 \cB_3 \les \int_{\bR^d} \sup_{\rho \leq t \leq t'} \int_\rho^t \int_{\bR^d} \int_{\bR^d} G^{2p'}_{t-s}(x-y) \|D^2_{(r,z),(\theta,w)}u(s,y)\|^{2p'}_{p}  dx dy ds dz. \\
\]
Applying \eqref{int G p} with respect to $dx$,
\[
 \cB_3 \les \int_{\bR^d} \sup_{\rho \leq t\leq t'} \int_\rho^t \int_{\bR^d} \|D^2_{(r,z),(\theta,w)}u(s,y)\|^{2p'}_{p} dy ds dz \les \int_{\bR^d} \int_\rho^{t'} \int_{\bR^d} \|D^2_{(r,z),(\theta,w)}u(s,y)\|^{2p'}_{p} dy ds dz. \\
\]
Relabelling, 
\[
 \cB_3 \les \int_\rho^{t'} \int_{\bR^d} \int_{\bR^d} \|D^2_{(r,z),(\theta,w)}u(t,x)\|^{2p'}_{p} dx dz dt. \\
\]
Replacing the $dx$ integral (which depends on $t$) by its supremum over $s \in [\rho, t]$,
\[
 \cB_3 \les \int_\rho^{t'} \int_{\bR^d} \sup_{\rho \leq s \leq t} \int_{\bR^d} \|D^2_{(r,z),(\theta,w)}u(s,x)\|^{2p'}_{p} dx dz dt. 
\]
For $0 \leq t \leq t' \leq T$ we have $1 \leq T^{1-p'}(t'-t)^{p'-1}$. We will applying this trivial bound so that $\cB_3$ matches the estimate for $\cB_5$ below, 
\[
 \cB_3 \les \int_\rho^{t'} (t'-t)^{p'-1} \int_{\bR^d} \sup_{\rho<s<t} \int_{\bR^2} \|D^2_{(r,z),(\theta,w)}u(s,x)\|^{2p'}_{p} dx dz dt.
\]

For $\cB_4$ we use Lemma \ref{Hol-lem} with respect to the measure $\mathbf 1_{\{|x-z| < T\}}dx$, if $p'< 1$. (If $p'  = 1$, the relation below is trivial. We obtain: 
\begin{align*}
 \cB_4 & \les \int_{\bR^d} \sup_{\rho \leq t \leq t'}  \bigg(\int_\rho^t \int_{\bR^d} \int_{\bR^d} \int_{\bR^d} G_{t-s}(x-y)G_{s-r}(y-z)G_{s-\theta}(y-w) f(y-y'), \\
 & \quad G_{t-s}(x-y') G_{s-r}(y'-z)G_{s-\theta}(y'-w) dx dy dy' ds \bigg)^{p'} dz.
 \end{align*}
Identifying the convolution $K_{t-s}(y-y')$, taking the supremum over $t \in [\rho,t']$, and using the fact that $t \to K_t(x)$ is non-decreasing, we have:
\begin{align*}
 \cB_4 & \les \int_{\bR^d} \sup_{\rho \leq t \leq t'}  \bigg(\int_\rho^t \int_{\bR^d} \int_{\bR^d} G_{s-r}(y-z)G_{s-\theta}(y-w) f(y-y') K_{t-s}(y-y') \\
 & \quad G_{s-r}(y'-z)G_{s-\theta}(y'-w)  dy dy' ds \bigg)^{p'} dz \\
 & \leq \int_{\bR^d} \bigg(\int_\rho^{t'} \int_{\bR^d} \int_{\bR^d} G_{s-r}(y-z)G_{s-\theta}(y-w) f(y-y') K_{t'}(y-y') \\
 & \quad G_{s-r}(y'-z)G_{s-\theta}(y'-w)  dy dy' ds \bigg)^{p'} dz.
\end{align*}
We apply again Lemma \ref{Hol-lem} with respect to the measure $\mathbf 1_{\{|z-w| < 2T\}}dz$, if $p' < 1$. (If $p'=1$, the relation below is trivial.) We obtain:
\begin{align*}
 \cB_4 & \les \bigg(\int_\rho^{t'} \int_{\bR^d} \int_{\bR^d} \int_{\bR^d} G_{s-r}(y-z)G_{s-\theta}(y-w) f(y-y')  K_{t'}(y-y') \\
 & \quad G_{s-r}(y'-z)G_{s-\theta}(y'-w) dz dy dy' ds \bigg)^{p'}.
\end{align*}
Identifying the convolution $K_{s-r}(y-y')$, and bounding it by $K_{t'}(y-y')$, we obtain:
\begin{align*}
 \cB_4 & \les \bigg(\int_\rho^{t'} \int_{\bR^d} \int_{\bR^d} G_{s-\theta}(y-w) G_{s-\theta}(y'-w) f(y-y') K_{t'}^2(y-y')  dy dy' ds \bigg)^{p'}.
\end{align*}
By Lemma \ref{fK int},% and the definition of $p'$, 
\begin{align*}
 \cB_4 & \les \bigg(\int_\rho^{t'} \bigg( \int_{\bR^d} G^{2p'}_{s-\theta}(y-w) dy \bigg)^{\frac{1}{p'}} ds\bigg)^{p'}.
\end{align*}
Using \eqref{int G p} then integrating with respect to $ds$,
\begin{align*}
 \cB_4 & \les \bigg(\int_\rho^{T} (s-\theta)^{\frac{2-2p'}{p'}} ds\bigg)^{p'} \les C_{T}.
\end{align*}

For $\cB_5$, we note that Corollary \ref{D2u-support} introduces the indicator $\mathbf 1_{\{|x-z| < T\}}$. We apply again Lemma \ref{Hol-lem} for the measure $\mathbf 1_{\{|x-z| < T\}}dx$, if $p'< 1$. (If $p'=1$, the relation below is trivial.) We obtain:
\begin{align*}
  \cB_5 & \les \int_{\bR^d} \sup_{\rho \leq t \leq t'}  \bigg(\int_\rho^t  \int_{\bR^d} \int_{\bR^d} \int_{\bR^d} G_{t-s}(x-y) \|D^2_{(r,z),(\theta,w)}u(s,y)\|_{p} f(y-y') \\
 & \quad G_{t-s}(x-y') \|D^2_{(r,z),(\theta,w)}u(s,y')\|_{p} dx dy dy' ds \bigg)^{p'} dz \\
 & \leq \int_{\bR^d} \bigg(\int_\rho^{t'}  \int_{\bR^d} \int_{\bR^d}  \|D^2_{(r,z),(\theta,w)}u(s,y)\|_{p} f(y-y')K_{t'}(y-y')  \\
 & \quad \|D^2_{(r,z),(\theta,w)}u(s,y')\|_{p} dy dy' ds \bigg)^{p'} dz.
\end{align*}
Applying Lemma \ref{fK int}, we get: % and the definition of $p'$,
\begin{align*}
  \cB_5 & \les \int_{\bR^d} \bigg(\int_\rho^{t'}  \bigg(\int_{\bR^d}  \|D^2_{(r,z),(\theta,w)}u(s,y)\|^{2p'}_{p} dy \bigg)^{\frac{1}{p'}} ds \bigg)^{p'} dz.
\end{align*}
The $dy$ integral depends on $s$. Taking the supremum of similar integrals for $s' \in [\rho,s]$, we obtain:
\begin{align*}
  \cB_5 & \les \int_{\bR^d} \bigg( \int_\rho^{t'}  \bigg(\sup_{\rho \leq s'\leq s} \int_{\bR^2}  \|D^2_{(r,z),(\theta,w)}u(s',y)\|^{2p'}_{p} dy \bigg)^{\frac{1}{p'}} ds \bigg)^{p'} dz. 
\end{align*}


We define $$M_{r,z,\theta,w}(t) = \sup_{\rho \leq s \leq t} \int_{\bR^d}  \|D^2_{(r,z),(\theta,w)}u(s,y)\|^{2p'}_{p} dy.$$ 

Then $M_{r,z,\theta,w}(t)$ is non-decreasing in $t$. Therefore, by Lemma \ref{incr ineq}, we have: 
\begin{align*}
  \cB_5 & \les  \int_{\bR^d} p'\int_\rho^{t'} (t'-s)^{p'-1} M_{r,z,\theta,w}(s) ds dz \\
  & \les \int_\rho^{t'} (t'-t)^{p'-1} \int_{\bR^d} \sup_{\rho \leq s \leq t} \int_{\bR^d}  \|D^2_{(r,z),(\theta,w)}u(s,y)\|^{2p'}_{p} dy dz dt.
\end{align*}

Combining the bounds gives the recursion, 
\begin{align*}
 & \int_{\bR^d} \sup_{\rho < t < t'} \int_{\bR^d} \|D^2_{(r,z),(\theta,w)}u(t,x)\|_{p}^{2p'} dx dz \\
 & \quad \leq A_T + B_T\int_\rho^{t'} (t'-t)^{p'-1} \int_{\bR^d} \sup_{\rho<s<t} \int_{\bR^d}  \|D^2_{(r,z),(\theta,w)}u(s,y)\|^{2p'}_{p} dy dz dt
\end{align*}
Let 
\[
 H_{r,\theta,w}(t) = \int_{\bR^d} \sup_{\rho<s<t} \int_{\bR^d}  \|D^2_{(r,z),(\theta,w)}u(s,y)\|^{2p'}_{p} dy dz.  
\]
Applying the above recursion twice and integrating yields, 
\[
H_{r,\theta,w}(t') \leq A_T + B_T \int_\rho^{t'} (t'-t)^{2p'-1} H_{r,\theta,w}(t)dt \leq   A_T + B_T \int_\rho^{t'} H_{r,\theta,w}(t)dt. 
\]
Applying the classical Gronwall lemma gives the desired bound, 
\[
 H_{r,\theta,w}(t) \leq A_Te^{B_T(t-\rho)} \leq C_T.
\]
\end{proof}

\begin{remark}
In the case $d = 1$ it is possible to establish the bound 
\[
 \|D^2_{(r,z),(\theta,w)} u(t,x)\|_p \les G_{t-r}(x-z)G_{2t-r-\theta}(z-w)
\]
 for all $0<\theta<r<t<T$ and $x,z,w \in \bR$, which is stronger than Theorem \ref{key-D2u-int}. However, this relies on the wave kernel $G_t$ being in $L^2(\bR^d)$ which only holds in dimension $d = 1$. The method is similar to the proof of Theorem \ref{key-Du}. We omit the proof for brevity. 
\end{remark}



%% END Key Estimates
%% BEGIN Body
\section{Proof of the main results}

In this section, we present the proofs of Theorems \ref{ergodic-th}, \ref{cov-th}, \ref{QCLT} and \ref{FCLT}.

%% BEGIN Ergodicity
\subsection{Ergodicity}

The proof of Theorem \ref{ergodic-th} follows directly from Theorem \ref{key-Du} and the more general result given below.

\begin{lemma}\label{ergodicity}
Let $\{Z(t,x) : t \geq 0,\ x \in \bR^d \}$ be an adapted random field such that $\{ Z(t,x)\}_{x \in \bR^d}$ is strictly stationary, and $Z(t,x) \in \bD^{1,2}$ for any $t \geq 0$ and $x \in \bR^d$. Assume that for any $0 < r < t$ and $x,z \in \bR^d$,
\begin{equation}\label{DZ}
\| D_{r,z}Z(t,x)\|_2 \leq C_t G_{t-r}(x-z)
\end{equation}
where $C_t > 0$ is a constant that depends on $t$. Then $\{ Z(t,x)\}_{x \in \bR^d}$ is ergodic.
\end{lemma}
\begin{proof} By Lemma 4.2 of \cite{BZ24}, it suffices to show that for any $k\geq 1$, and $b_1, \ldots , b_k, \zeta_1, \ldots, \zeta_k\in\bR^d$ and $g \in \{x \to \sin(x), x \to \cos(x)\}$, the following relation holds:
\begin{align}
\label{ERG1}
\lim_{R\to\infty} \frac{1}{R^{2d}}
{\rm Var}\bigg(  \int_{[0,R]^d} g \bigg(\sum_{j=1}^k b_j  Z(t,x + \zeta_j)  \bigg) dx \bigg)
=0.
\end{align}
We denote 
\[
X_{b,\zeta}(t,x) := g \left(\sum_{j=1}^k b_j  Z(t,x + \zeta_j)  \right)
\]
and 
\[
 V_R := {\rm Var}\left(  \int_{[0,R]^d} X_{b,\zeta}(t,x) dx \right)
\]
By Poincar\'e inequality \eqref{Poincare}, \eqref{LHS-q'}, and Minkowski's inequality,
\begin{align*}
V_R & \leq \int_{\bR_+} \int_{\bR^{2d}} \left\| D_{r,y}\int_{[0,R]^d} X_{b,\zeta}(t,x) dx  \right\|_2 \left\| D_{r,y'}\int_{[0,R]^d} X_{b,\zeta}(t,x) dx  \right\|_2 f(y-y') dydy'dr \\
 & \les \int_{\bR_+} \left[\int_{\bR^{d}} \left\| \int_{[0,R]^d} D_{r,y}X_{b,\zeta}(t,x) dx  \right\|_2^{2q'} dy\right]^{1/q'} dr \\
 & \les \int_{\bR_+} \left[\int_{\bR^{d}} \left( \int_{[0,R]^d} \left\| D_{r,y}X_{b,\zeta}(t,x)  \right\|_2 dx\right) ^{2q'} dy\right]^{1/q'} dr.
\end{align*}
By the definition of $X_{b,\zeta}(t,x)$, the chain rule \eqref{chain}, and \eqref{DZ},
\[
 \left\| D_{r,y}X_{b,\zeta}(t,x)  \right\|_2 \leq C_t\sum_{j=1}^k |b_j|  G_t(x + \zeta_j - y).
\]
Therefore, by the above and H\"older's inequality, 
\begin{align*}
V_R & \les  \int_{\bR_+} \left[\int_{\bR^{d}} \left( \int_{[0,R]^d} \sum_{j=1}^k |b_j| G_{t-r}(x + \zeta_j - y) dx\right) ^{2q'} dy\right]^{1/q'} dr \\
 & \les \sum_{j=1}^k |b_j|^2 \int_{\bR_+} \left[\int_{\bR^{d}} \left( \int_{[0,R]^d} G_{t-r}(x + \zeta_j - y) dx\right) ^{2q'} dy\right]^{1/q'} dr.
\end{align*}
For $R \geq t$, since $|x| \leq \sqrt{d}R$, we have $|\zeta_j - y| \leq (1+\sqrt{d})R$ for each $j = 1,\cdots,k$. Therefore,
\[
V_R \les \sum_{j=1}^k |b_j|^2 \int_{\bR_+} \left[\int_{B_{(1+\sqrt{d})R}(\zeta_j)} \left( \int_{[0,R]^d} G_{t-r}(x + \zeta_j - y) dx\right) ^{2q'} dy\right]^{1/q'} dr.
\]
Bounding $[0,R]^d$ by $\bR^d$ and integrating in the order of $dx$, $dy$ then $dr$ yields, 
\[
V_R \les \sum_{j=1}^k |b_j|^2 R^{\gamma}.
\]
Since $\gamma < 2d$ this proves Lemma \ref{ergodicity} and in turn, Theorem \ref{ergodic-th}. 
\end{proof}


%% END Ergodicity
%% BEGIN Limiting Cov
\subsection{Limiting covariance}

Let $\omega_d={\rm Vol}(B_1)$, i.e. $\omega_1=2$ and $\omega_2=\pi$.

\medskip

{\em Proof of Theorem \ref{cov-th}:} {\em Case 1. $f \in L^1(\bR^d)$.} By Lemma 18 of \cite{dalang99}, the process $\{u(t,x);t\geq 0,x\in \bR^d\}$ satisfies property (S) of \cite{dalang99}. In particular, this implies that for any $(t_1,x_1),\ldots,(t_k,x_k) \in  \bR_{+} \times \bR^d$ and $z \in \bR^d$,
\[
\big(u(t_1,x_1+z),\ldots,u(t_k,x_k+z)\big) \stackrel{d}{=} \big(u(t_1,x_1),\ldots,u(t_k,x_k)\big),
\]
where $\stackrel{d}{=}$ denotes equality in distribution. In particular, for any $t_1,t_2>0$,
\[
{\rm Cov}\big(u(t_1,x),u(t_2,y)\big)=:\rho_{t_1,t_2}(x-y) \quad \mbox{depends only on $x-y$}.
\]
By the dominated convergence theorem,
\begin{align*}
\frac{1}{R^{d}}\bE[F_R(t_1)F_R(t_2)] & =\frac{1}{R^{d}}\int_{B_R} \int_{B_R} \rho_{t_1,t_2}(x-y)dxdy=\frac{1}{R^{d}} \int_{B_{2R}} \rho_{t_1,t_2}(z) {\rm Vol} \big(B_R \cap B_R(-z)\big)dz \\
& \longrightarrow \omega_{d} \int_{\bR^d} \rho_{t_1,t_2}(z) dz=:K(t_1,t_2), \quad \mbox{as} \ R \to \infty,
\end{align*}
provided that 
\begin{equation}
\label{rho-int}
\int_{\bR^d} | \rho_{t_1,t_2}(x)| dx<\infty.
\end{equation}


To prove \eqref{rho-int}, we use Poincar\'e inequality \eqref{Poincare}, followed by the key estimate given by Theorem \ref{key-Du} and relation \eqref{int-G}:

\begin{align*}
& \int_{\bR^d}| \rho_{t_1,t_2}(x)| dx \leq \int_0^t \int_{\bR^d} \int_{\bR^d} \int_{\bR^d}  \|D_{r,z}u(t,x)\|_2 \|D_{r,z'}u(t,0)\|_2 f(z-z') dzdz'dx dr \\
& \leq C_{t,2}^2 \int_0^t \int_{\bR^d} \int_{\bR^d} \int_{\bR^d} G_{t-r}(x-z) G_{t-r}(-z') f(z-z') dzdz'dx dr\\
& = C_{t,2}^2\|f\|_{L^1(\bR^d)} \int_0^t (t-r)^2 dr<\infty.
\end{align*}


{\em Case 2. $f(x)=|x|^{-\beta}$ for some $\beta \in (0,d\wedge 2)$.}
Note that $F_R(t)=A_R(t)+B_R(t)$, where 
\begin{align*}
A_R(t) & =\int_0^t \int_{\bR^d} \varphi_{t,R}(s,y) \sigma \big( u(s,y)\big) W(ds,dy)\\
B_R(t) & = \int_0^t \int_{\bR^d} \varphi_{t,R}(s,y)\Big(b\big( u(s,y)\big)-\bE[b\big(u(s,y)\big)] \Big) dsdy.
\end{align*}
Hence, $\bE[F_{R}(t_1)F_R(t_2)]=\sum_{i=1}^4 T_{i,R}$ where
\begin{align*}
T_{1,R} &=\bE[A_{R}(t_1) A_{R}(t_2)], \quad T_{2,R}=\bE[A_R(t_1) B_R(t_2)] \\
T_{3,R} & =\bE[A_R(t_2) B_R(t_1)], \quad  T_{4,R}=\bE[B_R(t_1) B_R(t_2)].
\end{align*}

We will show that the limiting behaviour is dictated by $T_{1,R}$ and the other terms are negligible.
We examine each term separately. By It\^o isometry and the stationarity of $\{u(t,x)\}_{x\in \bR^d}$,
\[
T_{1,R}= \int_0^{t_1 \wedge t_2} \int_{(\bR^d)^2} \varphi_{t_1,R}(s,y) \varphi_{t_2,R} (s,z) |y-z|^{-\beta}\bE\big[\sigma(u(s,y)) \sigma(u(s,z))\big] dydz ds.
\]
Hence, $R^{\beta-2d} T_{1,R}=\cI_R+\cJ_R$, where
\begin{align*}
\cI_{R} &=R^{\beta-2d}\int_0^{t_1 \wedge t_2} \int_{(\bR^d)^2} \varphi_{t_1,R}(s,y) \varphi_{t_2,R} (s,z) |y-z|^{-\beta}{\rm Cov}\Big(\sigma(u(s,y)), \sigma(u(s,z))\Big) dydz ds, \\
\cJ_R &=\int_0^{t_1\wedge t_2}\xi^2(s) \psi_{R}(t_1,t_2;s,s) ds,
\end{align*}
where $\xi(s) = \bE[\sigma(u(s,0))]$, and for any $0\leq s_1 \leq t_1$ and $0\leq s_2 \leq t_2$, we denote
\[
\psi_R(t_1,t_2;s_1,s_2):=R^{\beta-2d} \int_{\bR^d} \int_{\bR^d} \varphi_{t_1,R}(s_1,y) \varphi_{t_2,R}(s_2,z)|y-z|^{-\beta}dydz.
\]
Similarly to Lemma 2.2 of \cite{BNZ}, it can be proved that
$\psi_R(t_1,t_2;s_1,s_2)$ is bounded by $c_{d,\beta}K_{d,\beta} (t_1-s_1)(t_2-s_2)$, and
\begin{equation}
\label{lem22}
\lim_{R \to \infty}\psi_R(t_1,t_2;s_1,s_2) =c_{d,\beta}K_{d,\beta} (t_1-s_1)(t_2-s_2)
\end{equation}
where, $K_{d,\beta}=\int_{\bR^d}|\cF \mathbf 1_{B_1}(\xi)|^2 |\xi|^{-(d-\beta)} d\xi$.

Using the same argument as in the proof of Proposition 3.1 of \cite{BNZ}, from the Clark-Ocone formula and the key estimate given by Theorem \ref{key-Du}, we deduce that
\[
\lim_{|y-z| \to \infty}{\rm Cov}\Big(\sigma(u(s,y)), \sigma(u(s,z))\Big) =0, 
\]
uniformly in $s\in [0,t_1 \wedge t_2]$, and therefore, $\cI_R \to 0$ as $R \to \infty$. 
By relation \eqref{lem22} and the dominated convergence theorem, 
\[
\lim_{R \to \infty}\cJ_R=c_{d,\beta} K_{d,\beta}\int_0^{t_1\wedge t_2}\xi^2(s) (t_1-s)(t_2-s)  ds=:K(t_1,t_2).
\]

Using the same idea, we study $T_{4,R}$:
\[
T_{4,R}= \int_0^{t_1} \int_0^{t_2} \int_{(\bR^d)^2} \varphi_{t_1,R}(s_1,y) \varphi_{t_2,R} (s_2,z) {\rm Cov}\Big(b(u(s_1,y)), b(u(s_2,z))\Big) dydz ds_1 ds_2.
\]
As in the proof of (3.3) of \cite{BNZ}, to show that $R^{\beta-2d}T_{4,R}\to 0$ as $R \to \infty$, it is enough to prove that
\begin{equation}
\label{cov-b}
\lim_{|y-z|\to \infty}{\rm Cov}\Big(b(u(s_1,y)), b(u(s_2,z))\Big) = 0,
\end{equation}
uniformly in $s_1\in [0,t_1]$ and $s_2 \in [0,t_2]$. 
To prove this, we use Clark-Ocone formula, followed by Jensen's inequality for conditional expectation, Theorem \ref{key-Du}, and boundedness of $b'$. We infer that:
\begin{align*}
& \Big|{\rm Cov}\Big(b(u(s_1,y)), b(u(s_2,z))\Big)\Big| \\
& \quad\leq \int_{0}^{s_1 \wedge s_2} \int_{(\bR^d)^2} 
\big\|D_{r,w} b\big(u(s_1,y)\big)\big\|_2 \big\|D_{r,w'} b\big(u(s_2,z)\big)\big\|_2 |w-w'|^{-\beta}dw dw' dr\\
& \quad \les \int_{0}^{s_1 \wedge s_2} \int_{(\bR^d)^2} G_{s_1-r}(w-y) G_{s_2-r}(w'-z)|w-w'|^{-\beta} dw dw' dr \\
& \quad \les \Big(|y-z|-(t_1+t_2)\Big)^{-\beta} \int_{0}^{s_1 \wedge s_2} \int_{(\bR^d)^2} G_{s_1-r}(w-y) G_{s_2-r}(w'-z)dw dw' dr .
\end{align*}
For the last line, we used the fact that $G_{s_1-r}(w-y) G_{s_2-r}(w'-z)$ contains the indicator of the set $\{|w-y|<s_1-r,|w'-z|<s_2-r\}$, and on this set
\[
|w-w'|\geq |y-z|-|(y-w)-(z-w')| \geq |y-z|-(s_1+s_2)\geq |y-z|-(t_1+t_2).
\]
By \eqref{int-G}, the last integral is equal to
$\int_0^{s_1 \wedge s_2} (s_1-r)(s_2-r)dr$. This proves that relation \eqref{cov-b} holds, uniformly in $s_1 \in [0,t_1]$ and $s_2 \in [0,t_2]$.

Finally, the fact that $R^{\beta-2d}T_{i,R}\to 0$ for $R \to \infty$ and $i=2,3$ follows by Cauchy-Schwarz inequality, using arguments similar to those used for $T_{1,R}$ and $T_{4,R}$ with $t_1=t_2$.
This concludes the proof of Theorem \ref{cov-th}.

%% END Limiting cov
%% BEGIN QCLT
\subsection{Quantitative CLT}

In this section, we give the proof of Theorem \ref{QCLT}. We recall that $|\cH|$ is the set of all measurable functions $\varphi: \bR_{+}\times \bR^d \to \bR$ such that
\[
\int_{\bR_{+}} \int_{(\bR^d)^2}|\varphi(t,x) \varphi(t,x')| f(x-x') dxdx'dt<\infty,
\]
and $|\cH^{\otimes 2}|$ is the set of all measurable functions $\varphi: (\bR_{+}\times \bR^d)^2 \to \bR$ such that
\[
\int_{\bR_{+}^2} \int_{(\bR^d)^4}|\varphi(t,x,s,y) \varphi(t,x',s,y')| f(x-x') f(y-y')dxdx'dydy'dt ds<\infty.
\]

We will use the following result, which is a particular case of Proposition 1.8 of \cite{BNQZ} with $\gamma_0=\delta_0$.

\begin{proposition}
\label{prop18}
   If $F\in\mathbb{D}^{2,4}$  has mean zero and variance $\sigma^2\in(0,\infty)$ such that with probability 1,  $DF\in| \cH|$ and $D^2F\in|\cH^{\otimes 2}|$, then
\[
d_{\rm TV}\Big( \frac{F}{\sigma},  Z\Big) \leq \frac{4}{\sigma^2} \sqrt{\mathcal{A}},
\]
where $Z\sim N(0,1)$ and
\begin{align*}
\mathcal{A}:&= \int_{\bR_+^3\times\bR^{6d}} \| D_{(r,z),(\theta,w)}^2 F \|_4  \| D_{(s,y),(\theta,w')}^2 F \|_4  \| D_{r,z'}F\|_4 \| D_{s, y' }F \|_4  \\
&\quad  \quad \quad f(y-y')f(z-z') f(w-w')  dydy'dzdz' dwdw'  dr ds d\theta.
\end{align*}
\end{proposition}

\bigskip



{\em Proof of Theorem \ref{QCLT}:} In the case $R < t$ the total variation is bounded by $1$ so $d_{TV}\left(\frac{F_R(t)}{\sigma_R(t)},Z\right) \leq t^{-(\gamma-2d)/2}R^{(\gamma-2d)/2}$. Hence, we may assume that $R \geq t$. By Proposition \ref{prop18},
\[
d_{\rm TV}\Big( \frac{F_R(t)}{\sigma_R(t)},  Z\Big) \leq \frac{4}{\sigma_R^2(t)} \sqrt{\mathcal{A}_R},
\]
where $Z\sim N(0,1)$ and
\begin{align*}
\mathcal{A}_R:&= \int_{[0,t]^3 \times\bR^{6d}}  \| D_{(r,z),(\theta,w)}^2 F_R(t) \|_4  \| D_{(s,y),(\theta,w')}^2 F_R(t)\|_4  \| D_{r,z'}F_R(t)\|_4 \| D_{s, y' }F_R(t) \|_4  \\
&\quad  \quad \quad f(y-y') f(z-z') f(w-w') dydy'dzdz' dwdw'  dr ds d\theta.
\end{align*}
By Theorem \ref{cov-th}, $\sigma_R^2(t)\sim C R^{\gamma}$. Therefore, it is enough to prove that 
\begin{equation}
\label{bound-AR}
\cA_R \les R^{3\gamma-2d}
\end{equation}
By Minkowski's inequality,
\begin{align*}
\mathcal{A}_R & \leq  \int_{[0,t]^3 \times\bR^{6d}} \int_{B_R^4} \| D_{(r,z),(\theta,w)}^2 u(t,x_1) \|_4  \| D_{(s,y),(\theta,w')}^2 u(t,x_2)\|_4  \| D_{r,z'}u(t,x_3)\|_4 \| D_{s, y' }u(t,x_4) \|_4  \\
&\quad  \quad \quad f(y-y') f(z-z') f(w-w') d\pmb{x}_4 dydy'dzdz' dwdw'  dr ds d\theta \\
\end{align*}
where $\pmb{x}_4=(x_1,x_2,x_3,x_4)$. We apply Theorem \ref{key-Du},
\begin{align*}
\cA_{R} & \les \int_{[0,t]^3 \times\bR^{6d}} \int_{B_R^4} \| D_{(r,z),(\theta,w)}^2 u(t,x_1) \|_4  \| D_{(s,y),(\theta,w')}^2 u(t,x_2)\|_4 \\
& \quad  \quad G_{t-r}(x_3-z') G_{t-s}(x_4-y') f(y-y') f(z-z') f(w-w') d\pmb{x}_4 dydy'dzdz' dwdw'  dr ds d\theta.
\end{align*}
By the definition of $\varphi$ given in \eqref{def-var}, we write $\cA_{R}$ as, 
\begin{align*}
\cA_{R} & \les \int_{[0,t]^3 \times\bR^{2d}} f(w-w')  \\
& \qquad \int_{\bR^d} \int_{B_R}\| D_{(r,z),(\theta,w)}^2 u(t,x_1) \|_4 dx_1 \int_{\bR^d} \varphi_{t,R}(r,z') f(z-z') dz' dz \\
& \qquad \int_{\bR^d} \int_{B_R} \| D_{(s,y),(\theta,w')}^2 u(t,x_2)\|_4  dx_2 \int_{\bR^d} \varphi_{t,R}(s,y') f(y-y') dy' dy \\
& \qquad  dwdw'  dr ds d\theta.
\end{align*}
By Corollary \ref{D2u-support}, $\| D_{(r,z),(\theta,w)}^2 u(t,x_1) \|_4$ has support $|x_1-z| < t$. Therefore, the integrand with respect to $dz$ is supported on $B_{R+t}$. Moreover, for $R \geq t$ we have $B_{R+t} \subseteq B_{2R}$. The same logic for the integral with respect to $dy$ yields, 
\begin{align*}
\cA_{R} & \les \int_{[0,t]^3 \times\bR^{2d}} f(w-w')  \\
& \qquad \int_{B_{2R}} \int_{B_R}\| D_{(r,z),(\theta,w)}^2 u(t,x_1) \|_4 dx_1 \int_{\bR^d} \varphi_{t,R}(r,z') f(z-z') dz' dz \\
& \qquad \int_{B_{2R}} \int_{B_R} \| D_{(s,y),(\theta,w')}^2 u(t,x_2)\|_4  dx_2 \int_{\bR^d} \varphi_{t,R}(s,y') f(y-y') dy' dy \\
& \qquad  dwdw'  dr ds d\theta.
\end{align*}

Since $0 \leq \varphi_{t,R}(r,z') \leq t-r$ we have that, 
\begin{align*}
\sup_{z \in B_{2R}} \int_{B_{2R}}\varphi_{t,R}(r,z')f(z-z')dz' \leq (t-r)\int_{B_{4R}}f(z')dz' \les (t-r)R^{\gamma-d}.
\end{align*}
Applying this twice to $A_{R}$ gives,
\begin{align*}
\cA_{R} & \les R^{2\gamma-2d}\int_{[0,t]^3} (t-r)(t-s) \int_{\bR^d}\int_{\bR^d}f(w-w') \\
& \qquad \left(\int_{B_{2R}} \int_{B_R} \| D_{(r,z),(\theta,w)}^2 u(t,x_1) \|_4 dx_1 dz\right) \\
& \qquad \left(\int_{B_{2R}} \int_{B_R} \| D_{(s,y),(\theta,w')}^2 u(t,x_2)\|_4 dx_2 dy\right)\\
& \qquad dwdw'drdsd\theta.
\end{align*}
By Corollary \ref{D2u-support}, $\| D_{(r,z),(\theta,w)}^2 u(t,x_1) \|_4$ has support $|x_1-w| < t$. Therefore, the integrand with respect to $dw$ is supported on $B_{R+t}$. Moreover, for $R \geq t$ we have $B_{R+t} \subseteq B_{2R}$. The same logic for the integral with respect to $dw'$ yields, 
\begin{align*}
\cA_{R} & \les R^{2\gamma-2d}\int_{[0,t]^3} (t-r)(t-s) \int_{B_{2R}}\int_{B_{2R}}f(w-w') \\
& \qquad \left(\int_{B_{2R}} \int_{B_R} \| D_{(r,z),(\theta,w)}^2 u(t,x_1) \|_4 dx_1 dz\right) \\
& \qquad \left(\int_{B_{2R}} \int_{B_R} \| D_{(s,y),(\theta,w')}^2 u(t,x_2)\|_4 dx_2 dy\right)\\
& \qquad dwdw'drdsd\theta.
\end{align*}
Applying \eqref{LHS-q'} and bounding $B_R, B_{2R}$ by $\bR^d$ gives, 
\begin{align*}
\cA_{R} & \les R^{2\gamma-2d}\int_{[0,t]^3}(t-r)(t-s) \\
& \quad \left\| \int_{\bR^d} \int_{\bR^d} \| D_{(r,z),(\theta,w)}^2 u(t,x_1) \|_4 dx_1 dz \right\|_{L^{2q'}(B_{2R})} \\
& \quad \left\| \int_{\bR^d}\int_{\bR^d} \| D_{(s,y),(\theta,w')}^2 u(t,x_2)\|_4 dx_2 dy \right\|_{L^{2q'}(B_{2R})} drdsd\theta.
\end{align*}
By Corollary \ref{D2u-support} $\| D_{(r,z),(\theta,w)}^2 u(t,x_1) \|_4$ is supported on $|x_1-w| < t$ and $|z-w| < 2t$. Let $p'$ be defined as in the statement of Theorem \ref{key-D2u-int}. Applying H\"older's inequality with power $2p'$ and measure $\mathbf 1_{|x_1-w| < t}\mathbf 1_{|z-w| < t} dx_1dz$,
\[
 \int_{\bR^d} \int_{\bR^d} \| D_{(r,z),(\theta,w)}^2 u(t,x_1) \|_4 dx_1 dz \les \left(\int_{\bR^d} \int_{\bR^d} \| D_{(r,z),(\theta,w)}^2 u(t,x_1) \|^{2p'}_4 dx_1 dz \right)^{1/2p'} 
\]
Taking the supremum,   
\begin{align*}
 \int_{\bR^d} \int_{\bR^d} \| D_{(r,z),(\theta,w)}^2 u(t,x_1) \|_4 dx_1 dz & \les \left(\int_{\bR^d} \sup_{\rho < s < t}\int_{\bR^d} \| D_{(r,z),(\theta,w)}^2 u(s,x_1) \|^{2p'}_4 dx_1 dz \right)^{1/2p'}, 
\end{align*}
where $\rho = \max(r,\theta)$. Applying Theorem \ref{key-D2u-int} gives,
\begin{align*}
 \int_{\bR^d} \int_{\bR^d} \| D_{(r,z),(\theta,w)}^2 u(t,x_1) \|_4 dx_1 dz & \les C_T 
\end{align*}
We have similar calculations for $\int_{\bR^d}\int_{\bR^d} \| D_{(s,y),(\theta,w')}^2 u(t,x_2)\|_4 dx_2 dy$. Therefore,   
\begin{align*}
\cA_{R,1} & \les R^{2\gamma-2d}\int_{[0,t]^3}(t-r)(t-s) \left(\int_{B_{2R}} C_T^{2q'} dw\right)^{1/2q'}  \left(\int_{B_{2R}} C_T^{2q'} dw'\right)^{1/2q'}  drdsd\theta. \\
& \les R^{3\gamma-2d} \int_{[0,t]^3} (t-r)(t-s) drdsd\theta \les R^{3\gamma-2d}.
\end{align*}
\qed
%% END QCLT
%% BEGIN FLCT
\subsection{Functional CLT}

In this section, we give the proof of Theorem \ref{FCLT}.

{\em Step 1. (finite-dimensional convergence)} In this step, we show that for any integer $m\geq 1$ and for any $t_1,\ldots,t_m>0$,
\[
\Big(\frac{1}{R^{\gamma/2}} F_R(t_1),\ldots, \frac{1}{R^{\gamma/2}}F_R(t_m)  \Big)
\stackrel{d}{\to} \Big(\cG(t_1),\ldots, \cG(t_m) \Big) \quad \mbox{as $R \to \infty$}.
\]
By Cram\'er-Wold theorem, this is equivalent to showing that for any $b_1,\ldots,b_m \in \bR$,
\[
X_R:=\frac{1}{R^{\gamma/2}}\sum_{j=1}^{m}b_jF_R(t_j) \stackrel{d}{\to}X:=\sum_{j=1}^{m}b_j \cG(t_j) 
\quad \mbox{as $R \to \infty$}.
\]
As in the proof of Theorem 1.5.(ii) of \cite{BS26}, it is enough to prove that $X_R/\tau_R \stackrel{d}{\to} Z$ as $R \to \infty$, where $\tau_R^2={\rm Var}(X_R) \to \tau^2={\rm Var}(X)$ and $Z \sim N(0,1)$. For this, we apply Proposition \ref{prop18} to $F=X_R$. It follows that
\begin{equation}
\label{bound-BR}
d_{\rm TV}\Big( \frac{X_R}{\tau_R},  Z\Big) \leq \frac{4}{\tau_R^2} \sqrt{\cB_R} \leq C \sqrt{\cB_R},
\end{equation}
where 
\begin{align*}
\cB_R:&= \int_{\bR_+^3\times\bR^{6d}} dr ds d\theta  dzdz' dydy' dwdw'  \\
&\quad  f(z-z') f(w-w') f(y-y')  \| D_{(r,z),(\theta,w)}^2 X_R \|_4  \| D_{(s,y),(\theta,w')}^2 X_R \|_4  \| D_{r,z'} X_R\|_4 \| D_{s, y' }X_R \|_4.
\end{align*}

Applying Minkowski's inequality to each of the four $\|\cdot\|_4$-norms above, it follows that
\[
\cB_R \leq \frac{1}{R^{2\gamma}} \sum_{i,j,k,\ell=1}^{m} \big|b_i b_j b_k b_{\ell}\big| \cA_R(t_i,t_j,t_k,t_{\ell}), 
\]
where
\begin{align*}
\cA_R(t_i,t_j,t_k,t_{\ell}) &= \int_{B_R^4} dx_1 dx_2 dx_3 dx_4 \int_{\bR_+^3 \times \bR^{6d}}  dr ds d\theta  dzdz' dydy' dwdw'  f(z-z') f(w-w') f(y-y')  \\
&\quad   \| D_{(r,z),(\theta,w)}^2 u(t_i,x_1) \|_4  \| D_{(s,y),(\theta,w')}^2 u(t_j,x_2) \|_4  \| D_{r,z'} u(t_k,x_3)\|_4 \| D_{s, y'} u(t_{\ell},x_4) \|_4.
\end{align*}
Using the same method as for the proof of \eqref{bound-AR}, it can be proved that
\[
\cA_R(t_i,t_j,t_k,t_{\ell}) \leq C R^{3\gamma-2d}
\]
Therefore, it follows that:
\[
\cB_R \leq CR^{\gamma-2d}
\]
Combining this with \eqref{bound-BR}, we infer that $d_{\rm TV}\Big( \frac{X_R}{\tau_R},  Z\Big) \to 0$ as $R \to \infty$.

\medskip

{\em Step 2. (tightness and H\"older continuity)}  We  will use Kolmogorov-Chentsov criterion (Theorem 3.23 of \cite{kallenberg02}) for the existence of the H\"older continuous modification, and the classical moment criterion (Corollary 16.9 of \cite{kallenberg02}) for tightness.  

Note that
\begin{align*}
u(t,x)-\bE[u(t,x)]& =\int_0^t \int_{\bR^d} G_{t-r}(x-z) \sigma\big(u(r,z)\big) W(dr,dz)+\\
&\int_0^t \int_{\bR^d} G_{t-r}(x-z) \Big(b\big(u(r,z)\big) -\bE\big[b\big(u(r,z)\big) \big]\Big) drdz.
\end{align*}

Integrating $dx$ on $B_R$, and using the definition of $\varphi_{t,R}$ \eqref{def-var} , we obtain that,
\begin{align*}
F_R(t) & =\int_0^t \int_{\bR^d} \varphi_{t,R}(r,z) \sigma\big(u(r,z)\big) W(dr,dz)+\\
&\int_0^t \int_{\bR^d}  \varphi_{t,R}(r,z)  \Big(b\big(u(r,z)\big) -\bE\big[b\big(u(r,z)\big) \big]\Big) drdz.
\end{align*}

Let $s,t \in [0,T]$ with $s<t$. We write $F_R(t)-F_R(s)=\sum_{i=1}^4 T_i$, where
\begin{align*}
T_1 & = \int_0^s \int_{\bR^d} \Big(\varphi_{t,R}(r,z) - \varphi_{s,R}(r,z) \Big) \sigma\big(u(r,z)\big) W(dr,dz)\\
T_2 &= \int_0^s \int_{\bR^d}  \Big( \varphi_{t,R}(r,z) - \varphi_{s,R}(r,z) \Big) \Big(b\big(u(r,z)\big) -\bE\big[b\big(u(r,z)\big) \big]\Big) drdz\\
T_3&=\int_s^t \int_{\bR^d} \varphi_{t,R}(r,z) \sigma\big(u(r,z)\big) W(dr,dz)\\
T_4&=\int_s^t \int_{\bR^d}  \varphi_{t,R}(r,z)  \Big(b\big(u(r,z)\big) -\bE\big[b\big(u(r,z)\big) \big]\Big) drdz.
\end{align*}
To estimate $T_1$ and $T_3$, we use Burkholder-Davis-Gundy inequality, followed by Minkowski inequality and the uniform moment estimate \eqref{def-K}. We obtain that for any $p\geq 2$,
\begin{align*}
\|T_1\|_p^2 & \les \int_0^s \int_{\bR^d} \int_{\bR^d}\Big(\varphi_{t,R}(r,y) - \varphi_{s,R}(r,y) \Big)  \Big(\varphi_{t,R}(r,z) - \varphi_{s,R}(r,z) \Big) f(y-z)dydz dr \\
\|T_3\|_p^2 & \les \int_s^t \int_{\bR^d} \int_{\bR^d} \varphi_{t,R}(r,y) \varphi_{t,R}(r,z) f(y-z) dydzdr.
\end{align*}
 
Applying \eqref{LHS-q'}, \eqref{est} and \eqref{diff est} gives
\begin{align*}
\|T_1\|_p^2 & \les \int_0^s \| \varphi_{t,R}(r,\cdot) - \varphi_{s,R}(r,\cdot)  \|_{L^{2q'}(\bR^d)}^2 dr \les R^{\gamma}(t-s)^{2/d}\\
\|T_3\|_p^2 & \les \int_s^t \| \varphi_{t,R}(r,\cdot) \|_{L^{2q'}(\bR^d)}^2 dr \les R^{\gamma}(t-s)^3 \les  R^{\gamma}(t-s)^{2/d}.
\end{align*}
For $T_2$ and $T_4$, we use Clark-Ocone formula, noting that $T_2$ is $\cF_s$-measurable, $T_4$ is $\cF_t$-measurable, and $\bE[T_2] = \bE[T_4] = 0$,
\begin{align*}
T_2 & = \int_0^s \int_{\bR^d} \bE[D_{r',y} T_2 |\cF_{r'}] W(dr',dy), \\
T_4 & = \int_0^t \int_{\bR^d}\bE[D_{r',y} T_4  | \cF_{r'}] W(dr',dy). 
\end{align*}

Similarly, by Burkholder-Davis-Gundy and Minkowski inequalities, for any $p\geq 2$,
\begin{align}
\label{T2}
\|T_2\|^2_p & \les \int_0^s \int_{\bR^d} \int_{\bR^d} \|D_{r',y} T_2\|_{p} \|D_{r',w} T_2\|_{p} f(y-w) dy dw dr', \\
\label{T4}
\|T_4\|^2_p & \les \int_0^t \int_{\bR^d} \int_{\bR^d} \|D_{r',y} T_4\|_{p} \|D_{r',w} T_4\|_{p} f(y-w) dy dw dr'. 
\end{align}
We now calculate $\|D_{r',y} T_2\|_{p}$ and $\|D_{r',y} T_4\|_{p}$. Using the chain rule \eqref{chain} for the Malliavin derivative, and relation \eqref{Mal-zero}, we have:
\begin{align*}
D_{r',y}T_2 &= \int_{r'}^s \int_{\bR^d}  \Big( \varphi_{t,R}(r,z) - \varphi_{s,R}(r,z) \Big) b'\big(u(r,z)\big) D_{r',y}u(r,z) drdz\\
D_{r',y}T_4&=\int_{s \vee r'}^t   \int_{\bR^d}  \varphi_{t,R}(r,z)  b'\big(u(r,z)\big) D_{r',y}u(r,z) drdz.
\end{align*}
Applying Minkowski's inequality, the fact that $b'$ is bounded, and Theorem \ref{key-Du}, we obtain that:
\begin{align*}
\|D_{r',y}T_2\|_{p} &\leq \int_{r'}^s \int_{\bR^d}  \left| \varphi_{t,R}(r,z) - \varphi_{s,R}(r,z) \right| \ \|D_{r',y}u(r,z)\|_{p}  drdz\\
&\les \int_{r'}^s  \int_{\bR^d}  \left| \varphi_{t,R}(r,z) - \varphi_{s,R}(r,z) \right| G_{r-r'}(z-y)  drdz\\
\|D_{r',y}T_4\|_p &\leq  \int_{s \vee r'}^t  \int_{\bR^d}  \varphi_{t,R}(r,z) \|D_{r',y}u(r,z)\|_{p} drdz \\
 &\les \int_{s \vee r'}^t \int_{\bR^d}  \varphi_{t,R}(r,z) G_{r-r'}(z-y) drdz. 
\end{align*}
Substituting the above bounds into \eqref{T2}, \eqref{T4} and applying \eqref{LHS-q'} gives, 
\begin{align*}
\|T_2\|^2_p & \les \int_0^s \left(\int_{\bR^d} \left( \int_{r'}^s \int_{\bR^d}  \left| \varphi_{t,R}(r,z) - \varphi_{s,R}(r,z) \right| G_{r-r'}(z-y)  drdz \right)^{2q'}  dy \right)^{1/q'}dr' \\
\|T_4\|^2_p & \les \int_0^t \left(\int_{\bR^d} \left( \int_{s \vee r'}^t \int_{\bR^d}  \varphi_{t,R}(r,z) G_{r-r'}(z-y)  drdz \right)^{2q'}  dy \right)^{1/q'}dr'.
\end{align*}
Note that $2q' > 1$ and $1/q' > 1$ so we can apply H\"older's inequality twice to the $dr$ integral, bounding the mass by $T$, to obtain, 
\begin{align*}
\|T_2\|^2_p & \les \int_0^s \int_{r'}^s \left(\int_{\bR^d} \left( \int_{\bR^d}  \left| \varphi_{t,R}(r,z) - \varphi_{s,R}(r,z) \right| G_{r-r'}(z-y)  dz \right)^{2q'}  dy \right)^{1/q'}drdr' \\
\|T_4\|^2_p & \les \int_s^t \int_{0}^r \left(\int_{\bR^d} \left(  \int_{\bR^d}   \varphi_{t,R}(r,z) G_{r-r'}(z-y)  dz \right)^{2q'}  dy \right)^{1/q'}dr'dr.
\end{align*}
By Young's inequality, 
\begin{align*}
\|T_2\|^2_p & \les \int_0^s \int_{r'}^s  \| \varphi_{t,R}(r,\cdot) - \varphi_{s,R}(r,\cdot) \|_{L^{2q'}(\bR^d)}^2 \|G_{r-r'}\|_{L^1(\bR^d)}^2 drdr' \\
\|T_4\|^2_p & \les \int_s^t \int_0^r  \| \varphi_{t,R}(r,\cdot)\|_{L^{2q'}(\bR^d)}^2 \|G_{r-r'}\|_{L^1(\bR^d)}^2 dr'dr.
\end{align*}
Finally, we apply \eqref{est}, \eqref{diff est} and $\|G_{r-r'}\|_{L^1(\bR^d)} = r-r'$, to obtain the estimate:
\begin{align*}
\|T_2\|^2_p & \les \int_0^s \int_{r'}^s  R^{\gamma} (t-s)^{\frac{2}{d}} (r-r')^2 drdr' \les R^{\gamma}(t-s)^{2/d} \\
\|T_4\|^2_p & \les \int_s^t \int_0^r R^{\gamma} (t-r)^{2} (r-r')^2 dr'dr \les R^{\gamma}(t-s)^{3} \les R^{\gamma}(t-s)^{2/d}.
\end{align*}
Hence, for any $p\geq 2$, $0\leq s<t\leq T$ and $R>T$,
\[
\|F_R(t)-F_R(s)\|_p \les R^{\gamma/2} (t-s)^{1/d}.
\]

%% END FCLT

%% END Body
%% BEGIN Appendix
\appendix
\section{Some auxiliary results}

In the proof of Theorem \ref{key-D2u-int}, we use the following application of H\"older's inequality, which allows us to move the exponent $p' \in (0,1)$ {\em outside} the integral.


\begin{lemma}
\label{Hol-lem}
Let $(X,\mathcal{X},\mu)$ be a finite measure space and $\phi:X \to \bR_{+}$ be a measurable function. For any $p' \in (0,1)$,
\[
\int_{X} \big(\phi(x)\big)^{p'} \mu(dx) \leq \big[\mu(X)\big]^{1-p'} \left( \int_{X} \phi(x) \mu(dx)\right)^{p'}.
\]
\end{lemma}

\begin{proof}
Let $p=\frac{1}{p'}>1$ and $F(x)=\big(\phi(x)\big)^{p'}$. By H\"older's inequality,
\[
\left(\int_{X} F(x) \mu(dx)\right)^p \leq \big[\mu(X)\big]^{p-1}  \int_{X} \big(F(x) \big)^p \mu(dx).
\]
The conclusion follows taking power $1/p$.
\end{proof}

In the proof of Theroem \ref{key-D2u-int}, we also use the following, which allows us to move the exponent $p' \in (\frac12,1]$ {\em inside} the integral, at a cost. 

\begin{lemma} \label{incr ineq}
Let $p' \in (\frac12,1]$ and $0 \leq r<t$. Let $M:[r,t] \to [0,\infty)$ be a non-decreasing right-continuous function. Then
\[
 \left( \int_r^t M(s)^{\frac{1}{p'}} ds \right)^{p'} \leq p' \int_r^t (t-s)^{p'-1} M(s) ds.
\]
\end{lemma}

\begin{proof} 
The $p'= 1$ case is trivial.
We suppose that $p'\in (\frac12,1)$.  By a change of variable,
\[
\left( \int_r^t M(s)^{\frac{1}{p'}} ds \right)^{p'} =\left( \int_0^{t-r} M(t-s)^{\frac{1}{p'}} ds \right)^{p'}.% \quad \int_r^t (t-s)^{p'-1} M(s) ds=\int_0^{t-r} s^{p'-1} M(t-s)ds.
\]
We denote $H(s)=M(t-s)$ for $s \in [0,t-r]$ and $L=t-r$. Then $H$ is non-increasing and left-continuous. It is enough to prove that
\begin{equation}
\label{H-ineq}
J:=\left( \int_0^L H(s)^{\frac{1}{p'}} ds \right)^{p'} \leq p' \int_0^L s^{p'-1} H(s) ds.
\end{equation}

By Minkowski's inequality,
\begin{align*}
J  &= \left[ \int_0^L \left( \int_0^{\infty} 1_{\{r\leq H(s)\}} dr\right)^{\frac{1}{p'}}ds \right]^{p'} = \left\| \int_0^{\infty}1_{\{r \leq H(s)\}} dr\right\|_{L^{\frac{1}{p'}}([0,L];ds)}\\
& \leq \int_0^{\infty} \|1_{\{r \leq H(s)\}} \|_{L^{\frac{1}{p'}}([0,L];ds)} dr=\int_0^{\infty} \|1_{\{s \leq H^{-1}(r)\}} \|_{L^{\frac{1}{p'}}([0,L];ds)} dr=\int_0^{\infty}\big(H^{-1}(r)\big)^{p'}dr,
\end{align*}
where $H^{-1}(y)=\sup\{x \in [0,L];H(x)\geq y\}$. On the other hand, by Fubini theorem,
\begin{align*}
p' \int_0^L s^{p'-1} H(s) ds&=p' \int_0^L s^{p'-1} \left(\int_0^{\infty} 1_{\{r \leq H(s)\}} dr \right) ds\\
& = p' \int_0^{\infty} \int_0^L s^{p'-1} 1_{\{s \leq H^{-1}(r)\}} dsdr=p' \int_0^{\infty} \left(\int_0^{H^{-1}(r)} s^{p'-1}ds \right) dr\\
&=\int_0^{\infty} \big(H^{-1}(r)\big)^{p'}dr.
\end{align*}
Relation \eqref{H-ineq} follows.
\end{proof}
%% END Appendix
%% BEGIN Bib
\begin{thebibliography}{99}
\bibitem{BHWXY25} Balan, R.M., Huang, J., Wang, X., Xia, P. and Yuan, W. (2025). Gaussian fluc-
tuations for the wave equation under rough perturbations. {\em Stoch. Proc. Their Appl.} {\bf 182}, 104569.

\bibitem{BNQZ} Balan, R.M., Nualart, D., Quer-Sardanyons, L. and Zheng, G. (2022).
The hyperbolic Anderson model: moment estimates of the Malliavin derivatives and applications. {\em Stoch. PDEs: Anal. \& Comp.} {\bf 10}, 757-827.

\bibitem{BS26} Balan, R.M. and Salins, M. (2026). Gaussian fluctuations for the nonlinear stochastic
heat equation with drift. {\em Elctr. J. Probab.} To appear.

\bibitem{BY24} Balan, R.M. and Yuan, W. (2024). Hyperbolic Anderson model with time-independent rough noise: Gaussian fluctuations. {\em Electr. J. Probab.} {\bf 29}, paper no. 177, 52 pp.

\bibitem{BZ24} Balan, R.M. and Zheng, G. (2024). Hyperbolic Anderson model with L\'evy white noise: spatial ergodicty and fluctuations. {\em Trans. AMS} {\bf 377}, 4171-4221.

\bibitem{BNZ} Bola\~{n}os-Guerrero, R., Nualart, D. and Zheng, G. (2021). Averaging 2d stochastic wave equation. {\em Electr. J. Probab.}
{\bf 26}, paper no. 102, 1-32.

\bibitem{dalang99} Dalang, R. C. (1999). Extending martingale
measure stochastic integral with applications to spatially
homogenous s.p.d.e.'s. {\em Electr. J. Probab.} {\bf 4}, no. 6,
1-29. %Erratum in {\em Electr. J. Probab.} {\bf 4}, no. 6, 1-5.

\bibitem{DNZ20} Delgato-Vences, F., Nualart, D. and Zheng, G. (2020). A central limit theorem for the stochastic wave equation with fractional noise. {\em Ann. Inst. Henri Poincar\'e: Probab.
Stat.} {\bf 56}, 3020-3042. 

\bibitem{HNV20} Huang J., Nualart D. and Viitasaari L. (2020). A central limit theorem for the stochastic heat equation. {\em Stoch. Proc. Their Appl.} {\bf 130}, 7170-7184. 

\bibitem{HNVZ20} Huang, J., Nualart, D., Viitasaari L. and Zheng, G. (2020). Gaussian fluctuations for the stochastic heat equation with colored noise. {\em Stoch. PDE: Anal. Comp.} {\bf 8}, 402-421. 

\bibitem{kallenberg02} Kallenberg, O. (2002). {\em  Foundations of Modern Probability}. Second Edition. Springer, New York.   
    
\bibitem{nualart06} Nualart, D. (2006). {\em The Malliavin Calculus and Related Topics}. Second edition. Springer-Verlag, Berlin.

\bibitem{NZ22} Nualart, D. and Zheng, G. (2022).
 Central limit theorems for stochastic wave equations in dimensions one and two. {\em Stoch.  Partial Differ. Equ.  Anal. Comput.} {\bf 10}, 392-418.

\end{thebibliography}
%% END Bib


\end{document} 
